\quellseite{123}

\section{Summarium}

Ist erstlich ein Exempel, wie Joseph des Landes Brauch hält, mit seines Vaters Tod, V. 2. Daß er ihn aber ins Land Canaan führet, und befiehlt, wenn seine Nachkommen aus Egypten ziehen, daß sie seine Gebeine auch sollen mitführen, V. 25., geschieht alles anzuzeigen, daß er fest an Gottes Verheißung halte vom Land Canaan, auf daß seine Nachkommen an solcher Verheißung auch nicht zweifeln.

Darnach ist ein schön Exempel hie, daß Joseph seinen Brüdern nicht allein von Herzen vergibt, und sie nicht entgelten läßt, sondern deutet solche Mißhandlung zum besten, V. 19., auf daß auch wir lernen Böses mit Gutem bezahlen, und von Herzen einer dem andern vergeben.

\section{Votum}

GOTT regiere unsere Herzen, daß wir nicht allein unserm Nächsten herzlich vergeben, sondern auch das Böse mit Gutem belohnen lernen. Und solches um unseres Heilandes willen, Amen.

\part{Das zweite Buch Mose}

\chapter{Das 1. Capitel}

\section{Vorrede}

Gleich wie im ersten Buch Mose wiederholet und erkläret wird die Verheißung von dem gebenedeyten Weibesamen und Schlangentreter Christo unserm HErrn, und werden darin angehänget herrliche Verheißungen, auch schöne Exempel, wie die Patriarchen in wahrem Glauben an die Verheißung sich gehalten haben; also lehret das andere Buch Mose, wie Gott die heiligen zehn Gebote wiederholt hat auf dem Berg Sinai, und erzählet schreckliche Exempel, wie Gott die gestraft, welche seine Gebote muthwillig übertreten haben. So helfe nun der hochgelobte Sohn Gottes, daß wir das Buch Mose also anhören, auf daß wir uns alle daraus bessern, damit unser Leben Gott wohlgefalle, und unserm Nächsten zu Nutz diene, und daß wir zeitlich und ewig gesegnet werden, Amen. Darauf höret nun an das 1. Capitel.

\section{Der Kinder Israel Dienstbarkeit und Drangsal in Egypten}
\quellseite{124}



\section{Summarium}

Ist eine Historia, in welcher man fürnehmlich siehet, wie der Teufel Gottes Segen begehret zu hindern, und des Josephs Treu und Wohlthat vergessen, und seine Nachkommen vom Tyrannen Pharao greulich beschweret werden.

Auf daß man lerne, daß die rechte Kirche oder Gottes Volk in der Welt allweg sich haben leiden müssen. Aber daneben hilft Gott, daß, je mehr sie gedrückt werden, je mehr sie sich mehren und ausbreiten.

Am Ende aber ist ein fein Exempel von den Hebammen, welche dem tyrannischen Befehl Pharao nicht wollten gehorsam sein, und darum von Gott gesegnet werden, V. 17. Denn in unbilligen Sachen ist man der Obrigkeit keinen Gehorsam schuldig, sondern man muß auf Gottes Gehorsam fürnehmlich sehen, und der Menschen Gehorsam dahinten lassen; nach dem Spruch Petri, Apost. 5, 29.: Man soll Gott mehr gehorchen denn den Leuten.

\section{Votum}

Der allmächtige und barmherzige GOtt erhalte seine Kirche, und thue wieder Gutes zeitlich und ewig allen, die seiner Gemeine und bedrängten Christen Gutes erzeigen, hochgelobt in Ewigkeit, Amen.

\chapter{Das 2. Capitel}

\section{Vorrede}

Wie kein edlerer Schatz auf dieser Welt zu finden ist, weder Gottes Wort, als darin allein die seligmachende Weisheit, Erkenntniß und Furcht Gottes, lebendiger Trost und ewiges Heil verborgen ist: also mögen und sollen sich die Leute wohl für die seligsten achten und halten, denen Gott solchen theuren Schatz vertrauet und gönnet, und dagegen für die unseligsten, die denselben nicht haben oder nicht haben wollen. So danket Gott alle herzlich für die besondere, ja für die höchste Gnade, daß ihr auch habt die heilige Biblia, und daß euch die Schriften der heiligen Propheten und Apostel zu Lehr, Trost und Warnung fürgelesen werden; auch wendet nöthigen Fleiß an, daß zu eurem großen Nutz solches geschehe, in Kraft des heiligen Geistes. Und höret jetzt an das 2. Capitel.

\section{Moses Geburt, Auferziehung, Flucht und Heirath}
\quellseite{125}

a Masa heißt ziehen; daher heißt Mose gezogen, nämlich aus dem Wasser.

b Gersom heißt ein Fremdling oder Ausländer.

†Non est in Hebræo.

‡(Elieser) heißt Gott mein Helfer.

\section{Summarium}

Die Historia von Mose hebt sich hie an, in welcher wir sehen, wie Gott der Welt und des Teufels Anschläge wunderbarlich brechen und zu nicht machen kann; weil eben des Königes Tochter Mosen erhält, V. 6., welcher des jüdischen Volks Herzog sein, und Egypten ganz wüst machen sollte.

Darum deutet das 11. Capitel B. 23. an die Ebräer, diese Historia also, und spricht:

Durch den Glauben ward Moses, da er geboren war, drei Monat verborgen von seinen Eltern, darum, daß sie sahen, daß er ein schön Kind war, und fürchteten sich nicht für des Königs Gebot. Durch den Glauben wollte auch Moses, da er groß ward, nicht mehr ein Sohn heißen der Tochter Pharao, und erwählete viel lieber mit dem Volk Gottes Ungemach zu leiden, denn die zeitliche Ergötzung der Sünde zu haben. Und achtete die Schmach Christi für größern Reichthum denn die Schätze Egyptens, denn er sahe an die Belohnung.

\section{Datum}

Gott unser Vater, der Mose erhalten und regieret hat, der wolle auch uns mit Gnaden bedenken, und wolle heut auch der Welt und des Teufels Anschläge hindern und brechen, Amen.
\quellseite{126}

\chapter{Das 3. Kapitel}

\section{Vorrede}

Die heilige Biblia ist das beste Lehrbuch und Trostbuch. Ein Lehrbuch ist die heilige Biblia: denn durch Gottes Wort und dessen Uebung und Betrachtung zündet Gott an und erwecket rechte Erkenntniß Gottes und wahren Glauben, unterrichtet auch, wie man Gott in rechtem Glauben und gutem Gewissen wahrhaftig und gefällig dienen, desgleichen wie ein Christ sein Leben und Stand recht anstellen und führen soll. Ein Trostbuch ist auch die heilige Biblia: denn daraus angefochtene, bekümmerte, gefängstigte Herzen zu jederzeit rechten, gewissen und beständigen Trost schöpfen und erlangen mögen, durch Hülfe des heiligen Geistes, daß Buch die Biblia ist, als des rechten himmlischen Kanzlers. So höret an das 3. Kapitel.

\section[Mose wird berufen, die Kinder Israel aus Egypten]{Mose wird berufen, die Kinder Israel aus Egypten zu führen}

a Ich werde seyn wenn ihr dahin kommet, so will ich bei euch seyn, und mich so erzeigen, daß ihr erkennen sollt, daß ichs sei.
\quellseite{127}

\section{Summarium}

Ist ein tröstlich Exempel, wie Gott dem Tyrannen nicht ewig zusehen, sondern seinem Volk, das unbillig beschweret ward, helfen und ihr Gebet erhören will, um seiner Verheißung willen. Darum offenbaret er sich Mose im brennenden Busch, B. 2., und fordert ihn zu solchem Werk und sagt, sein Name heiße Gott des Abraham, Gott des Isaak, Gott des Jakob, B. 6. Das ist, ein Gott, der sich durch seine Verheißung mit diesen Vätern verbunden habe, und ihnen zugesagt, daß er ihre Nachkommen nicht allein aus Egypten erlösen, und ins gelobte Land bringen wolle, sondern auch aus ihren Nachkommen den Samen kommen lassen, durch welchen alle Völker auf Erden sollen gesegnet, das ist, von Sünde und Fluch ledig, und ewig selig werden.

\section{Votum}

Der starke allmächtige GOtt erhalte auch bei uns und in der ganzen Welt seine Kirche; und verleihe, daß wir deren gläubige und gehorsame Gliedmaßen sein und allezeit bleiben, Amen.

\chapter{Das 4. Kapitel}

\section{Vorrede}

Von Natur haben wir allesamt erkaltete Herzen, die zu Gottes Wort keine Lust haben. Fabeln und Mährlein höret man viel lieber, da sperret man Maul und Nasen auf und ist niemand keine Zeit zu lang, man höret solches mit Lust und Liebe. In der Kirche aber wartet man nicht gerne lange. O man höret sich Gottes Wort und der Predigen bald überdrüssig. Dies betrachtet, und weil Gottes Wort eine Werkstatt ist des heiligen Geistes und der Wagen, darauf der heilige Geist zu uns kommt und bei uns eingehet, zu unserer Besserung und Trost, so sollt ihr nimmermehr Gottes Wortes überdrüssig werden. Und derowegen höret auch jetzt an das 4. Kapitel.

\section[Mose wird in seinem Beruf mit der Gabe]{Mose wird in seinem Beruf mit der Gabe Wunder zu thun gestärket}
\quellseite{128}

\section{Summarium}

Dies ist eine historia, wie Moses sich gewehret und dem Beruf nicht folgen wollte, V. 1. Darum, daß er wohl sahe, wie mancherley Gefahr es auf ihm hatte, so er wider des Königs Willen sollte die Juden aus dem Lande führen. Aber Gott tröstet ihn, er wolle bei ihm sein und ihn lehren, was er sagen soll, V. 3. Gibt ihm auch Macht, drei Wunderwerke zu thun, mit der Schlange, ausfätzigen Hand, und dem Blut; auf daß er ein gewiß Zeugniß hätte, was er thäte, daß ers aus Gottes Befehl thäte. Applicatio. Solches Exempel ist tröstlich allen denen, so Gottes Beruf vor ihnen haben, daß sie unverzagt sein und Gott sollen sorgen lassen, wie sie es ausrichten, was er ihnen befohlen hat. Gleich wie man siehet, wo der Beruf nicht ist, es sei Macht und Vermögen so groß es wolle, so kann man doch nichts Fürnehmliches ausrichten.

*Blutbräutigam) das ist, sie ward zornig, und sprach: Es kostet Blut, daß du mein Mann bist und muß mein Kind beschneiden; welches sie ungern that, als das eine Schande war unter den Heiden. Bedeutet aber des Gesetzes Volk, welches gerne wollte Gott haben, aber es will das Kreuz nicht leiden, noch den alten Adam beschneiden lassen, bis es thun muß.

\section{Summarium}

Ist erstlich ein tröstlich Exempel, wie Gott den Tyrannen wegnimmt, der Mose nach dem Leben stund, V. 19. Darnach, wie Moses, da er seinem Beruf nachkommen wollte, in die Gefahr kommt, daß ihn Gott tödten wollte, V. 23., welches darum geschah, eben wie oben mit Jacob, daß er über Gottes Verheißung etwas leiden, und mit dem Glauben darüber fechten müßte, auf daß er hernach desto fester stünde, und stärker würde. Denn unversuchte Leute taugen zum Predigtamt eben so wenig, als zu andern Händeln.
\quellseite{129}

\section{Votum}

GOtt der Vater des Lichts, von dem alle gute Gaben von oben herab kommen, segne unsern Beruf, unsere Aemter und alle christliche Händel und Orber, und wende ab alle Hindernisse und Gefahr, Amen.

\chapter{Das 5. Kapitel}

\section{Vorrede}

Gott strafet die Gottlosen und hilft den Frommen. Er gibt aber den Gottlosen eine zeitlang Glück und Wohlfahrt, ehe er ihnen nach ihrem Verdienst im Zorn ablohnet, und die Frommen läßt er das Kreuz auch eine zeitlang wohl drücken und fühlen, ehe denn er ihnen nach seiner Verheißung in Gnaden aushilft. Derowegen sollen wir uns in Gottes Weise schicken lernen, sollen Buße thun, geduldig sein, beten und Gottes Hülfe hoffen und erharren. So wird es denn zuletzt wohl gehen denen, die fromm bleiben und sich recht halten: aber die Gottlosen zuletzt ausgerottet und wie das Gras abgehauen, und wie das grüne Kraut verwelken werden, Psalm 37. Ferner höret an das 5. Kapitel.

\section[Mose wird von Pharao verachtet, und das Volk]{Mose wird von Pharao verachtet, und das Volk noch mehr geängstiget}

a Das ist, deine armen Leute müssen unrecht haben und Sünder sein, man suchet Schuld zu deinem Volk.
\quellseite{130}

\section{Summarium}

Sind zwei sonderliche Exempel: das erste, wie die Tyrannen und Gottes Feinde sich stellen, wenn ihr Verderben am nächsten ist, nämlich, daß sie von Gott nichts wollen wissen, V. 2., und mit den Leuten je länger je unbarmherziger umgehen, auf daß sie ihr eigen Urtheil forttreiben und zum Ende bringen, V. 6.

Das andere Exempel ist, daß wir an den Kindern Israel sollen lernen, wie sich Gott dazu pflege zu schicken, wenn er will helfen, nämlich, daß ers hinter sich anfühet, und wenn man der Hülfe gewartet, allererst am tiefsten ins Unglück fällt. Darum gilt es hie hoffen und harren, bis die Hülfe hernach kommt; wie es den Kindern Israel ging, ob es sich gleich im Anfang sehr übel dazu schickte.

\section{Votum}

Der treue GOtt, der ein GOtt ist alles Trostes und der Geduld, wirke und erhalte Geduld allen, die jetzt das liebe Kreuz tragen; hindere auch und stürze alle boshaftigen und unwilligen Feinde seines heiligen Namens, Amen.

\chapter{Das 6. Kapitel}

\section{Vorrede}

Wir müssen bekennen, daß wir alles Kreuz mit unserm sündhaftigen Wesen und Leben wohl verdienet haben. Wir sollen bekennen, daß es eine Strafe ist von Gott, viel geringer, denn unsre Sünden sind, Judith 8. Da bekennen sollen wir, daß wir auch die größten zeitliche, und darnach auch die ewige Strafe verdienet hätten, wenn Gott mit uns ins Gerichte gehen und mit uns handeln wollte nach seiner Gerechtigkeit und unserm Verdienst: wie wir denn solches alle bekennen müssen, wir haben so heilig gelebt, als wir immer vermocht. Und dieweil Gott unser dort mit der ewigen wohlverdienten Strafe um seines lieben Sohnes willen verschonen will, sollen wir uns nicht weigern allhier ein wenig Kreuz zu leiden, sondern sollen mit dem Propheten Micha sagen, Mich. 7.: Ich will den Zorn des HErrn tragen, denn ich habe ihm gesündiget. Davon höret an das 6. Kapitel.

\section[Mose bekommt von Gott einen neuen Befehl, beschreibet]{Mose bekommt von Gott einen neuen Befehl, beschreibet sein Geburtsregister}

a Nicht offenbaret) die Patriarchen haben Gott wohl erkannt, aber eine solche öffentliche gemeine Predigt war zu der Zeit von Gott noch nicht aufgegangen, wie durch Mose und Christum geschehen ist.

b Meine Hand gehabet) das ist, geschworen, wie oben. 1 Mos. 22, 16.
\quellseite{131}


\section{Summarium}

Ist eine Historia, wie Gott Mosen und das Volk tröstet, auf daß sie fest an Gottes Zusagung halten, und durch die Beschwerung der Frohndienste sich nicht sollten dahin bewegen lassen, daß sie darum ihrem Gott und Mosi seinem Diener, ungehorsam würden, V. 6.

Aber wie der Text die sagt, V. 9., sie hörten ihn nicht für der Angst und Seufzen, und für harter Arbeit.

Also gehet es noch, daß man sich des Unglücks halben allerweg mehr bekümmert, denn man sich Gottes Gnade und Gottes Worts tröstet; so doch Gottes Wort billig sollte von uns am größten geachtet werden, als der höchste Schatz und Trost in allerlei Jammer und Anfechtung.

\section{Votum}

Der barmherzige Gott lasse ihm alle befohlen sein, die noch das Creuz tragen, und in Noth und Gefahr stecken, und helfe ihnen an Leib und Seele, daß sie ihn ewig mit uns dafür loben und preisen, Amen.

\chapter{Das 7. Capitel}

\section{Vorrede}

Wir sind in dieser Welt in des Teufels Reich, nicht anders denn wenn ein Pilgrim in eine Herberge käme, da er wüßte, daß sie alle im Hause Räuber wären. Wenn er dahin kommen würde, würde er sich dennoch rüsten und aufs beste als er könnte vorsehen, und nicht viel schlafen: also sind wir jetzund auf Erden, da der böse Geist ein Fürst ist, und hat der Menschen Herzen in seiner Gewalt und thut durch sie was er will. Ferner, wie der Teufel einen Hirsch umbringet, wenn er ihm auf den Kopf springet und ihm das Gehirn ausfrißt, oder greift ihn bei der Kehle, und beißet sie ihm entzwei: also ist auch der Teufel, wenn er einen Menschen in seine Gewalt bekommt, so kann man seiner nicht leichtlich los werden, er führet den Menschen aus einer Sünde in die andere, und thut ihm Schaden an Leib
\quellseite{132}
und Seele. Derwegen wollen wir wachen und beten, daß der Teufel unser ja nicht mächtig werde, und uns nicht führe in Sünde, Schand und Schaden. Dieß betrachten wir auch bei dem 7. Kapitel.

\section{Verwandlung des Wassers in Blut}

a Gott kann Pharao weder mit Plagen, noch mit Wunder, noch mit Dräuen bessern, er muß ihn lassen ersaufen im Meer.
\quellseite{133}

\section{Summarium}

Dies ist eine Historia, wie Pharao der Tyrann sich an kein Zeichen kehret, sondern in seinem Unglauben und Ungehorsam bleibet, V. 4. Solches aber muß dazu dienen, auf daß Gott ihn desto härter angreife, und Gottes Ehre und Kraft desto herrlicher geoffenbaret werde in ganz Egypten. Eben also gehet es mit den gottlosen Tyrannen noch: je mehr sie sich wider Gottes Wort und Willen setzen, je größer und tröstlicher ist’s, daß sie am Ende darüber zu Boden gehen; Gottes Wort aber bleibet ewiglich.

\section{Votum}

Der HErr Zebaoth, der starke GOtt, steure allen Tyrannen, und tröste und erhalte dagegen auch alle, die seinen heiligen Namen anrufen, bekennen, loben und preisen, Amen.

\chapter{Das 8. Kapitel}

\section{Vorrede}

Gottes Strafen sollen unsere Bußprediger sein; und wenn Gott züchtiget, so soll man nicht ärger, sondern frömmer werden. Denn wie St. Paulus sagt, Röm. 6.: Unser alter Mensch wird gekreuziget, daß der sündhaftige Leib aufhöre, daß wir hinfort der Sünde nicht dienen. Eben dies meinet auch St. Petrus, da er schreibt, 1 Epist. Kap. 4.: Wer am Fleisch leidet, der höret auf von Sünden, daß er hinfort, was noch hinterstelliger Zeit im Fleisch ist, nicht der Menschen Lüsten, sondern dem Willen Gottes lebe. Derowegen wer von Gott gestraft wird, der thue Buße und lasse sich den Teufel und seine Werkzeuge nicht verführen, sondern folge Gottes Wort in Gehorsam, so wird ihm folgen Gottes Segen zur Wohlfahrt. Darauf höret an davon das 8. Kapitel.

\section{Egypten wird mit Fröschen, Läusen und Ungeziefer geplaget}

\section{Summarium}

Ist ein Exempel, wie die Tyrannen verzagt werden, wenn Gottes Gericht über sie kommt, V. 8. Aber so bald sie Luft haben, und die Strafe nachläßt, wieder hinfallen und bleiben wie sie gewesen sind, V. 15.

Daß aber der Egypter Zauberer etliche Zeichen nachthun, und endlich an den Läusen zu Schanden werden, V. 18., ist ein Zeugniß, daß solche Teufels-Künstler mehr nicht können, denn Gott aus seinem Zorn ihnen selbst, oder andern zur Aergerniß und Strafe verhänget, aber sie
\quellseite{134}
müssen endlich an dem Geringsten zu Spott und Schande werden.

a Ungeziefer) das die Griechen heißen Cynonomia, sind alle bösen Würmer, so da Schaden thun im Felde, Raupen, Fliegen, Zweyfalter, Ameisen, Käfer, Brenner und dergleichen Geschmeiß, das Bäume und Gewächse verderbet.

\section{Summarium}

Ist ein Exempel, wie die Tyrannen in der Noth sich füßlich demüthigen, aber doch immer ihrem eigenen Willen und nicht dem Wort Gottes folgen; wie man an unsern Papisten auch sieht: da hilft keine Strafe, kein Ueberweisen, kein Vermahnen, sondern eben wie Pharao, bleiben sie verstockt und ungebessert, ob sie gleich selbst bekennen müssen, daß ihr Wesen mit Gottes Wort nicht übereinstimmt. Auf solche Pharaonische Sünde gehört auch ein Pharaonisch plötzliches Verderben.

\section{Votum}

Der treue GOtt behüte uns für Pharaonischer Verstockung, und regiere uns durch den heiligen Geist, daß wir unsern Willen brechen, und seinem heiligen Willen nach seinem Wort folgen, bis zum seligen Todesstündlein, Amen.

\chapter{Das 9. Kapitel}

\section{Vorrede}

Des alten Adams oder des alten Menschen Kreuzigung steht in erster Buße, daß man Reue habe über die Sünde und sich fleißig für derselben hüte, und also das Fleisch im Zaum halte. Das kömmt manchem wohl so sauer an als der Tod, sonderlich wenn man von böser alter Gewohnheit soll ablassen, seinem eigenen
\quellseite{135}
Sünde abbrechen und steuren, böser Gesellschaft sich verziehen, der Welt Ungunst auf sich laden, ein Gewinnlein fahren lassen, so fort an. Derowegen weil wir, als die wir Christum angehören, unser Fleisch kreuzigen sollen, samt den Lüsten und Begierden: so laßt uns Gott um Hülfe darzu anrufen; denn die Natur kann ihr selber nicht widerstreben noch abbrechen. Darauf höret auch an das 9. Kapitel.

\section{Pestilenz. Schwarze Blattern. Hagel}

\section{Summarium}

Ist ein Exempel, wie Gott mit Strafen nicht aufhören wolle, wo man von Sünden nicht ablassen und seinem Wort nicht folgen will. Wiederum, wo man folget, da will er schützen und behüten, ob es gleich sonst allenthalben übel zugeht; wie er seine Juden behütet, da sonst das ganze Egypten um der Sünde und Ungehorsams willen mußte leiden.
\quellseite{136}

\section{Summarium}

Ist ein sonderlich Exempel wie oben, daß die verstockte Herzen sich fälschlich demüthigen, als wollten sie fromm werden, begehren Gnade und Fürbitte. Aber es ist ihnen nicht länger Ernst, bis das Unglück fürüber ist. Ob wohl solches frommen Herzen wehe thut, sollen sie doch diesen Trost daraus nehmen, daß die Gottlosen mit solcher falscher Demuth ihnen selbst ihre Sache nur ärger machen, und zu größerer Strafe und Zorn Gott verursachen.

\section{Votum}

Gott unser lieber Vater helfe uns allen, daß wir rechtschaffene Buße thun, und in der Frömmigkeit beständig fortfahren, und schütze und erhalte in diesen gefährlichen Zeiten seine Christenheit und alle Gläubige, Amen.

\chapter{Das 10. Capitel}

\section{Vorrede}

Wie Gott mit Strafen nicht aufhöret nach den Dräuungen des Gesetzes, wo man von Sünden nicht ablässt, und seinem Wort zu einem christlichen Leben nicht folget: also wo man folget, und einen ehrbaren heiligen Wandel führet, da beschützet und behütet Gott nach den Verheißungen des Evangelii, ob es gleich sonst allenthalben übel zugehet. Derwegen wer nicht will gestraft sein, der thue Buße und lebe beständig in Gottesfurcht; und wen Gott erhalten soll unter künftigem Unglück, der beharre auf dem Wege der Gerechtigkeit. Denn, wie an die Ebräer am 2. gesagt wird: Eine jegliche Uebertretung und Ungehorsam hat empfangen seinen rechten Lohn. Also spricht David Psalm 37.: Laß vom Bösen und thue Gutes, und bleibe immerdar. Dies lehret auch das 10. Kapitel.

\section{Heuschrecken und Finsterniß}
\quellseite{137}
* Die Heuschrecken heißen hie nicht Hagab auf Hebräisch, wie an etlichen Oertern, sondern Arbe. Es sind aber vierfüßige fliegende Thiere, und rein zu essen, wie Hagab, Lev. 11, 12. aber uns unbekannt, ohne daß sie Heuschrecken gleich sind.
\quellseite{138}

\section{Summarium}

Dies ist eine Historia, wie über die andern Plagen Gott Egyptenland mit den Heuschrecken verderbet, B. 13., und darnach mit Finsterniß geschlagen habe, B. 21., auf daß man lerne, wo man Gottes Wort nicht folgen will, daß da kein Aufhören der Strafe zu hoffen sei. Darum will auch Gott, daß solcher Exempel bei Kindes Kindern soll gedacht werden, auf daß man beides lerne: sich für solchem Gott fürchten, der die so ihm ungehorsam sind, dermaßen strafen will; und auf solchen Gott trauen, der seinem Volk so wunderlich helfen kann.

Es ist auch hie ein Exempel, wie Pharao je länger je verstockter wird, und die Kinder Israel nicht allein nicht ziehen, sondern auch Moses will tödten lassen, B. 18.

\section{Votum}

Gott, unser aller lieber Vater, behüte uns für Pharaonischer Blindheit, und verleihe, daß wir deinem Wort gehorsamlich folgen, und daß wir die gegenwärtigen Strafen zu wahrer Buße gebrauchen, Amen:

\chapter{Das 11. Kapitel}

\section{Vorrede}

Was Gott den Unbußfertigen dräuet, das läßt er endlich wahrhaftig über sie ergehen; und da schonet er keines, wer er auch ist vor der Welt. Denn es ist kein Ansehen der Person für Gott, und bei ihm trägt niemand vor sein großes Geschlecht, ehrliche Freundschaft, ansehnliches Alter, hohe Weisheit, vielfältige Erfahrung, große Kunst, herrliche Aemter, mächtige Gewalt und dergleichen, was den Menschen in der Welt bei den Leuten ein Ansehen machen kann. Deswegen so wollen wir unsere Sünde erkennen, die Gott abbitten, davon abstehen, und ein Christen-Leben führen, daß wir nicht zeitlichen und ewigen Jammer über uns führen. Denn um Jesu Christi willen vergibt Gott den Bußfertigen alle Sünden und lindert ihnen aus Gnaden die wohlverdienten Strafen. Darauf höret auch an das 11. Kapitel.

\section{Auszug aus Egypten befohlen}

a (Großer Mann) das saget er darum, daß es Wunder ist, wie er nicht getödtet ist von den Egypten; sie haben sich müssen eines Aergern und Aufruhrs fürchten.

\section{Summarium}

Dies ist eine Historia, wie Moses Pharao dräuet mit der letzten Strafe, B. 4. Aber Pharao achtet solches Dräuen nicht, bis ers erfährt, B. 9. Wie denn die Gottlosen allwege nicht eher glauben, bis der Glaube ihnen in die Hände kommt, und sie es erfahren. Daß aber der HErr heißet silbernes Geschirr von den Egypten entlehnen, und dasselbe ihnen entwenden, B. 2., ist darum geschehen, auf daß den Kindern Israel ihre Arbeit belohnet würde, welche ihnen vom Könige mit unbilliger Beschwerung ward aufgelegt. Denn Gott will, daß die Regimente billiger Weise geführt, und alle Tyrannei verhütet, und sonderlich einem jeden seine Mühe und Arbeit belohnet werde. Wo solches weltliche Obrigkeit nicht will thun, muß es Gott selbst schaffen, aber durch ein hartes Mittel. Denn die Egyptier hätten die Juden leichter können bezahlen, denn daß sie ihnen alle ihre Schätze und Kleinode anhängen, und darnach Leib und Leben auch darüber verlieren, wie im rothen Meer geschehen ist.

Digitized by Google
\quellseite{139}

\section{Votum}

Helfe GOtt, daß wir mit den Gottlosen sein Wort und Träumungen ja nicht verachten, sondern denen Glauben geben, und seinem Wort zu täglicher Besserung mit allen Frommen folgen, Amen.

\chapter{Das 12. Capitel}

\section{Vorrede}

Gott ist den Unbußfertigen ein eifriger Gott, den Gläubigen aber ist er ein gütiger Gott. Dies lehret er selber im 2 Buch Mose im 20. Kapitel, da er sagt: Ich der HErr dein Gott, bin ein eifriger Gott, denen, die mich hassen; und thue Barmherzigkeit an viel Tausenden, die mich lieb haben und meine Gebote halten. Derowegen wer unter den mannigfaltigen Strafen, damit der gerechte Gott jetzt die Welt wegen der überhäuften großen Sünden, Schande und Lastern heimsucht, will Schutz, Trost und Hülfe haben, der bitte heute abermal, Gott wolle ihm gnädig sein und alle seine Sünde vergeben; er halte sich auch in wahrem Glauben an Jesum Christum, unsern HErrn, und lebe nach seinem Willen; also wird er sein und bleiben ein Verborgener des HErrn, und wird mit Freuden singen mit dem König und Propheten David: Der HErr decket mich in seiner Hütten zur bösen Zeit; er verbirget mich heimlich in seinem Gezelt: denn der HErr verbirget die Seinen heimlich bei ihm, und verdecket sie in seiner Hütten, Ps. 27. und 31. Davon höret an das 12. Kapitel.

\section[Stiftung des Osterlamms. Erwürgerung der Erstgeburt. Des Auszugs]{Stiftung des Osterlamms. Erwürgerung der Erstgeburt. Des Auszugs Anfang}

a Zusammen kommet) das ist, sie sollen predigen, das Benedikte und Gratias beten für das Passah und Erlösung aus Egypten. Darum gebraucht er das Wörtlein Kara, welches heißt predigen, lesen in der Versammlung.
\quellseite{140}

\section{Summarium}

Ist eine Historia, wie und aus was Ursache Gott den Juden das Osterfest habe eingesetzt.

Zum ersten, daß es soll sein ein ewig Gedächtniß der herrlichen Erlösung aus Egypten. Denn Gott will, daß man seiner Wohlthat nicht vergessen, sondern ihm dafür danken und derselbigen in allen Nöthen zum Trost brauchen soll.

Zum andern, daß solch Osterfest auch soll ein Fürbild sein der geistlichen Erlösung von Sünde und Tod, durch das Blut Christi, des unschuldigen Lämmleins, wie St. Paulus solches deutet 1 Cor. 5, 7.: Daß unser Osterlamm sei Christus, für unsere Sünde geschlachtet. Solches Osterlamm sollen wir täglich essen, das ist, durch den Glauben uns seines Todes und Verdienstes täglich trösten; und den Sauerteig ausfegen, das ist, der Sünden wehren, und ihr nicht folgen.
\quellseite{141}

b) Passah) heißt ein Gang, darum, daß der HErr durch Egyptenland des Nachts ging und schlug alle Erstgeburt todt. Bedeutet aber Christus Sterben und Auferstehen, damit er von dieser Welt gegangen ist, und in demselben Sünde, Tod und Teufel geschlagen und uns aus dem rechten Egypten geführet hat zum Vater; das ist unser Passah oder Ostern.

\section{Summarium}

Ist ein tröstlich Exempel, in welchem wir sehen, wie Gott die Seinen endlich errettet und der Tyrann Pharao mit seinem harten Sperren nichts anderes hat ausgerichtet, denn daß er sich und das ganze Land Egypten in alles Unglück gebracht, und jetzt nach solchem Schaden ohne Dank das thun muß, daß er zuvor, ehe einiger Schaden geschehen war, zu Dank nicht hat wollen thun. Also muß Gott die Tyrannen zum Gehorsam bringen, wenn sie sein Wort so gar nicht achten und sich für ihm nicht demüthigen wollen. Wer auf die Historien und tägliche Erfahrung siehet, wird solcher Exempel viel finden, daß es den großen Herren an ihren Anschlägen fehlet und Gottes Wille mit Macht ungewehret fortgehet.

\section{Gebet}

Der treue GOtt erbarme sich heut auch seiner verfolgten Christenheit, und steure und wehre allen bösen Anschlägen und Tyrannen, um JEsu Christi seines Sohns, unseres HErrn willen, Amen.

\chapter{Das 13. Kapitel}

\section{Vorrede}

Wohl denen, die an Gottes Wohlthaten gedenken und ihm dafür dankbar sein. Denn je öfter und fleißiger ein Mensch Gott danket für die empfangene Wohlthaten, je mehr ihm Gott ferner, so lange er lebet, Gutes erzeiget. Dieses lehret David, der in der Person Gottes, Ps. 50. also saget: Wer Dank opfert, der preiset mich, und das ist der Weg, daß ich ihm zeige das Heil Gottes, das ist, mit Danksagung bauet man unserm HErrn Gott den Weg und lockt ihn, daß er mit Wohlthun desto öfter und reichlicher wiederkomme. Derwegen lobsinget eurem Gott, und danket ihm täglich für seine Gnade und Barmherzigkeit, wie er solche euch und den Euern, an Leib, Seele und Nahrung, erzeiget hat. Darauf höret nun an das 13. Kapitel.

\section{Heiligung der Erstgeburt. Des Auszugs Fortgang}

a) Abib) ist der Monat, den wir April heißen; denn die Hebräer heben ihr Neujahr an nach der Natur, wenn die Ding wieder neu grünet und wächst und sich fühlet. Darum heißet er auch Mensis novorum, daß denn Alles neu wird.
\quellseite{142}

b Ungesäuert Brod) so hart wird der Sauerteig verboten, daß man ja das lautere Evangelium und Gottes Gnade, nicht unser Werk und Gesetz soll predigen nach der Auferstehung Christi. Wie St. Paulus 1 Cor. 5, 6. auch zeiget. Und ist solch Essen nichts Anderes, denn glauben an Christum.

c Schilfmeer) die Griechen heißen es das rothe Meer, von dem rothen Sand und Boden; aber die Hebräer heißen es das Schilfmeer, von dem Schilf.

d Gerüstet) hebräisch gefünfft. Was das sei, lassen wir Andere suchen; ob es sei, daß sie bei fünf Haufen gezogen, oder bei fünf neben einander gegangen sind, oder was solch fünf sei. Denn von der fünften Rippe, so das Harnisch erreicht am Leibe (wie die Juden hie \unsicher{fabeln}), verstehen wir nichts.

\section{Summarium}

Ist ein neuer Befehl Gottes, wie man neben dem Opferfest alle Erstgeburt dem HErrn heiligen und mit einem Opfer lösen soll, V. 1. Zum Gedächtniß, daß Gott sein Volk aus Egypten habe ledig gemacht, da er alle ihre Erstgeburt durch den Engel in einer Nacht hat würgen lassen.

Solches ist uns zur Lehre fürgeschrieben, daß wir ja die Werke Gottes nicht aus der Acht lassen, sondern derselbigen gedenken, Gott dafür danken und uns damit in allerlei Anfechtung trösten sollen.

Das Feuer und Wolkensäule, in welcher Gott für seinem Volk herzog, V. 21., sind ein Fürbild des Wortes Gottes und des heiligen Geistes, dadurch wir hie auf Erden, gleich wie in einer Wüste geführt, getröstet und erhalten werden.

\section{Gebet}

GOtt der heilige Geist regiere uns durch sein Wort, daß wir alle Sünden fliehen, und alle Tugenden üben, und die Wallfahrt dieses Lebens in wahrem Glauben und gutem Gewissen seliglich vollenden, Amen.
\quellseite{143}

\chapter{Das 14. Kapitel}

\section{Vorrede}

Gott tröstet und erhält die Frommen und schrecket und stürzet die Gottlosen. Derowegen wer in Sünden lebet, der thue Buße, daß ihn Gott nicht strafen dürfe. Wer aber in Gottesfurcht lebet, und ist daneben ein betrübter Kreuzträger, oder stecket in Noth und Gefahr, der sei geduldig und bete; Gott wird ihn wahrhaftig erretten und erhalten. Diesen Trost lehret David im 93. Psalm und saget: Die Wasserwogen im Meer sind groß und brausen gräulich; der HErr ist noch größer in der Höhe; und St. Johannes, da er spricht 1 Epist. 4. Cap.: Der in euch ist, der ist größer, denn der in der Welt ist. Dies erinnert auch das 14. Kapitel.

\section[Des Auszugs Vollendung. Der Egypter Untergang im rothen]{Des Auszugs Vollendung. Der Egypter Untergang im rothen Meere}


a Heil) Hülfe.

b Schreiest du) merckt hie ein trefflich Exempel, wie der Glaube kämpfet, zappelt und schreyet in Nöthen und Fährlichkeit, und wie er sich an Gottes Wort bloß hält und von Gott Trost empfähet und überwindet.
\quellseite{144}

e (Erleuchtet) das ist, es war ein Wetterleuchten in den dicken Wolken.


\section{Summarium}

Ist eine treffliche Historia, in welcher wir sehen, wie Gott die Seinen, ehe er sie von den Egyptern gar erledigt, zuvor in einen neuen Schrecken und Gefahr kommen läßt, V. 5. Wie man sonst auch siehet, daß Gott nicht eher Errettung schaffet, es sei denn zuvor mit aller menschlichen Hülfe und Mittel verloren.

Da gehet des Glaubens Kampf an, daß man in solcher Noth allein auf Gottes Barmherzigkeit, Allmächtigkeit und Verheißung sich verlasse. Wie denn Moses seinen Juden auch hie prediget: Fürchtet euch nicht, und sehet zu, was der HErr heute für ein Wunderwerk an euch thun werde, V. 13.

Also kehret es sich denn um: das Volk Gottes ist am ersten erschrocken, und sieget doch endlich durch die Hülfe Gottes. Die Egypter sind am ersten sicher und freudig, endlich liegen sie im rothen Meer, da es am tiefsten ist. Dies soll das Ende sein aller Feinde Gottes und Tyrannen.

\section{Votum}

Wir danken Gott herzlich für beides, daß er an den Israeliten und Egyptern gethan hat; der strafe und stürze ferner alle Unbußfertige, und erhalte die Frommen in Gnaden, und segne und erhalte sie ewiglich, Amen.

\chapter{Das 15. Kapitel}

\section{Vorrede}

Nach dem Glauben ist die Danksagung der fürnehmste und vortrefflichste Gottesdienst. Derowegen wie der Vogel am Gesang und der Topf am Klang erkannt wird: also ist der Mensch ein rechter Christ, der gern betet und Gott für seine Wohlthaten herzlich danket. So wollen wir nun auf Gottes Barmherzigkeit und Gaben Achtung geben, und als wahre Christen wollen wir ihm auch täglich dafür danken. Und wollen also diesen Gottesdienst alhie anfangen, weil wir ihn fürnehmlich in jenem Leben mit den lieben Engeln Gottes treiben, und mit ihnen fröhlich und ewig singen werden: Heilig, heilig, heilig ist der HErr Zebaoth, Es. 6. Das lehret auch das 15. Kapitel.

\section{Lobgesang Mosis. Bitteres Wasser süß gemacht}
\quellseite{145}

a Mara) heißt bitter, und bedeutet Leiden und Anfechtung, welche durch das Kreuz Christi im Glauben auch süß werden, Matth. 11, 29. 30. Mein Joch ist süß.

\section{Summarium}

Ist ein Danklied für die herrliche Erlösung am rothen Meer, welche Moses rühmet für ein sonderliches Trostbild, daß sich Gott nicht allein mit großer Kraft wider seine Feinde leget, sondern auch mit rechter Gnade gegen seinem Volk hat sehen lassen. Auf daß auch wir auf solchen Gott unser Vertrauen setzen, und uns seiner Hülfe in allerlei Noth trösten lernen.

Aber wie bald sie solcher Wohlthat vergessen haben, zeigt an das Volk am dritten Tage hernach, welches, da es ihnen an Wasser mangelt, wider Mosen murrete, V. 24. Aber Gott hilft auch in derselbigen Noth, V. 25.; und verheißet durch Mosen ferner weiter zu helfen, V. 26., so fern man sich seines Worts und Befehls halten will.
\quellseite{146}
Amen.

Lob, Ehr und Preis sei GOtt für alle seine Wohlthaten; der regiere und befestige uns, daß wir ihm ferner täglich mit Mund, Herz und gottseligem Wandel danken bis zu dem ewigen Leben, Amen.

\chapter{Das 16. Kapitel}

\section{Vorrede}

Zwei Dinge unterscheiden uns fürnehmlich von den Heiden und Unchristen: als nämlich Gottes Wort und ein heiliges Leben. Beides bestätiget der königliche Prophet David, da er sagt im 93. Psalm: Dein Wort ist eine rechte Lehre, Heiligkeit ist die Zierde deines Hauses ewiglich. Dies Wort Gottes haben wir, und das wird gepredigt und vorgelesen, daß wir darnach recht gläuben, christlich leben und geduldig leiden sollen in Lieb und Leid, Noth und Tod. Derowegen wollen wir Gottes Wort also anhören, lernen und gebrauchen, daß wir gute Werke draus machen; wir wollen auch in allerlei löblichen Tugenden uns üben, nach Erheischung eines jeden Beruf und Amt, durch Hülfe und Regierung des heiligen Geistes, daß wir uns als wahre Christen gegen Gott und Menschen verhalten. Darauf höret ferner an das 16. Kapitel.

\section{Die Wachteln und das Manna wird gegeben}

a Man) heißt auf hebräisch eine Gabe. Bedeutet, daß uns das Evangelium ohne unser Verdienst und Zuthun, aus lauter Gnaden vom Himmel gegeben wird, wie das Man auch gegeben ward.
\quellseite{147}

\section{Summarium}

Sie gehet an der Juden Historia, darinn man siehet, wie ein ungläubig und störrig Volk sie gewesen sind. In Egypten schrieen sie sehr über die Arbeit und unbillige Beschwerung. Sie wünschen sie, daß sie wieder drinnen wären, um eines Bissen Fleisches und Stück Brodes willen, V. 2. Und denken nicht einmal daran, daß sie einen solchen Gott haben, der sich schon bewiesen habe, daß er in allerlei Nöthen könne und wolle helfen. Ja sie vergessen all seiner Wunder und sagen, Mose und Aaron haben sie auf nichts Gutes aus Egypten geführet, V. 3.

Aus dieser Historie sollen wir unsere Art und Natur erkennen lernen, wie wir so gar verderbet sind; aber Gott wiederum so barmherzig ist, daß er gerne vergibt, und Fleisch und Brod den ungläubigen murrenden Juden die Fülle gibt.


b Zeugniß das ist an dem Ort, da man opfert und betet, und der Predigtstuhl war, ehe die Hütte war gemacht.


\section{Summarium}

Es ist ein fein Exempel, wie der Unglaube nichts ausrichtet mit dem, da er die Sache verhoffet zu fördern. Gott wollte, wie wir auch im Vater Unser beten, man sollte am täglichen Brod sich genügen lassen. Aber etliche wollten fürsammlen, denen verdarb das Man, V. 18. Etliche wollten am Sabbath sammeln, die fanden nichts, V. 27. Auf daß wir beides lernen, am ersten nicht geizig sein, sondern am täglichen Segen uns genügen lassen, und für das Morgende nicht sorgen.

Zum andern, daß wir uns die Nahrung am rechten Gottesdienst nicht lassen hindern. Wie Christus auch befiehlt: Suchet am ersten das Reich Gottes, so soll euch das andere alles zufallen, Matth. 6, 33.
\quellseite{148}

\section{Votum}

Helfe der treue GOtt, daß wir fleißig thun was wir sollen, daß er auch als ein lieber Vater treulich thue, was wir zu seinen Ehren und unserer Nothdurft begehren und wollen, hochgelobet in alle Ewigkeit, Amen.

\chapter{Das 17. Kapitel}

\section{Vorrede}

Der weise Mann Sirach vermahnet ein jedes unter uns zur Danksagung und spricht im 32. Kapitel: Danke für alles dem, der dich geschaffen und mit seinen Gütern gesättiget hat. Und St. Paulus vermahnet jedermann zum Gebet und sagt Ephes. 6.: Betet stets in allem Anliegen, mit Beten und Flehen im Geist; aber, ach wir müssens alle bekennen, undankbar sind wir, wenn uns Gott geholfen hat, und nachlässig sind wir im Gebet, wenn wir Gottes Hülfe bedürfen. Derowegen laßt uns solche Unart erkennen und davon abstehen: lasset uns Gott für seine Wohlthaten danken; denn er verdient es um uns alle Augenblick. Lasset uns auch fleißig beten; denn zu dieser Zeit Gottes Ehre und unsere höchste Nothdurft solches erfordert. Davon unterrichtet auch das 17. Kapitel.

\section[Die Kinder Israel werden aus einem Felsen getränket,]{Die Kinder Israel werden aus einem Felsen getränket, überwinden die Amalekiter}




\section{Summarium}

Es ist ein Exempel, wie die Welt neben allem, was sie von Gott weiß und gesehen hat, alsobald sie einen Mangel spüret, unwillig wird, Gottes sich nicht tröstet und schlecht nichts leiden will; ja alle Schuld auf die Prediger und ihre Obrigkeit schiebet. Daß Moses sich hie bezeugen muß, sie möchten ihn mit Steinen zu Tode werfen, V. 4. Aber Gott beweiset hie wiederum, wie er mehr nach seiner Barmherzigkeit, denn nach dem wir werth sein, mit uns handeln wolle.
\quellseite{149}
Zum andern ist ein fein Exempel hie, mit Moses Gebet, daß man des Sieges von Gott gewarten und durchs Gebet Gott suchen soll, V. 11. Daß man aber Mose seinen Arm unterstützt, V. 12., bedeutet, wenn das Gebet soll stark und rechtschaffen gehen, daß man Buße thun und mit dem Leben sich bessern soll, daß beides beisammen stehe, ein ernstliches Gebet und christliches Leben.

\section{Votum}

Der barmherzige gütige GOtt helfe, daß wir in täglicher Buße leben und fleißig beten, daß GOtt uns auch beistehe, und allen unsern Feinden steure und wehre, auch uns unsere Nothdurft aus Gnaden beschere. Amen.

\chapter{Das 18. Kapitel}

\section{Vorrede}

Wo in einer Stadt oder Lande die Kirchen mit treuen Predigern und das weltliche Regiment mit treuer Obrigkeit begnadet ist, da wohnet Gott und da ist Glück, Heil und Seligkeit. Denn christliche Prediger sind nützliche Amtleute im geistlichen Reiche Gottes. Und fromme Obrigkeiten sind heilsame Amtleute im weltlichen Reiche Gottes. Diese Kleinode hat auch unsere Stadt und Vaterland. Derowegen laßt uns Gott dafür danken und Gottes Wort, daß unsere Seelsorger lehren, und der Obrigkeit, die uns mit göttlichen Gesetzen regieret, also gehorsam sein, daß gemeldete Gnadenschätze bis auf den jüngsten Tag uns erhalten werden, Amen. Dies unterrichtet auch das 18. Kapitel.

\section{Jethro gibt Mose einen guten Rath}
\quellseite{150}

\section{Summarium}

Ist eine Historia, wie Moses das Regiment mit mehr Leuten bestellet und er fürnehmlich des Predigens und Gottesdienstes wartet, V. 24. Denn es ist nicht eines Mannes Arbeit, einem solchen Haufen fürzustehen.

Es ist aber hie sonderlich fein abgemalet, was für Leute zum Regiment zu brauchen seien, V. 21. Nämlich redliche Leute, die Ehre und Tugend lieb haben, und fest darob halten. Darnach gottesfürchtige Leute, die wahrhaftig und nicht geizig sind. Denn wo Gottesfurcht nicht ist, da wird wenig Glück sein. Und wer nicht wahrhaftig noch treu ist, da ist man übel mit versorget in Händeln. Was aber geizige und eigennützige Leute sind, die taugen gar nichts zum Regiment. Denn der Eigennutz blendet sie, daß sie auf den gemeinen Nutz nicht können sehen.

\section{Votum}

Der barmherzige Gott erhalte und segne unsere liebe Obrigkeit, daß sie auf Erden also regiere, damit sie mit dem Himmel-Könige Christo, um seines Verdienstes willen, regiere in Ewigkeit, Amen.

\chapter{Das 19. Kapitel}

\section{Vorrede}

Daß Gott die heiligen Zehen Gebote auf dem Berge Sinai wiederholet hat, ist so eine große Wohlfahrt, daß sie keines Menschen Zunge genugsam ausreden kann.

Derowegen danken wir heute Gott für solche Wohlfahrt; der wolle uns allen den heiligen Geist geben, daß wir durch seine Regierung nach den heiligen Zehen Geboten leben, die darin verbotene Sünde fliehen und die darin befohlenen guten Werke thun und üben. Denn, wie die Kriegsleute einen Bestallungs-Brief haben, darauf sie schwören: also sind unsere Bestallungs-Briefe die heiligen Zehen Gebote, welche in den Schriften der Propheten und Apostel erkläret werden; und daß wir nach diesem Bestallungs-Brief leben wollen, haben wir alle in unserer Taufe geschworen, und derowegen wer dawider handelt, der wird an Gott der heiligen Dreifaltigkeit meineidig. Dieß betrachtet auch bei dem 19. Kapitel.

\section{Vorbereitung, das Gesetz Gottes anzuhören}
\quellseite{151}

\section{Summarium}

Ist eine Historia, wie GOtt sich geoffenbaret habe, daß er das Gesetz auf dem Berge Sinai gegeben hat. Und reimet sich solche Offenbarung fein mit der Predigt des Gesetzes. Denn gleich wie das Donnern und Blitzen, wo es so gräulich ist, erschreckt und furchtsam macht, v. 16.: also wird auch ein menschlich Herz, wo es des Gesetzes Predigt recht trifft, erschrocken seiner Sünden und Gottes Zorn halben. Darum heißt es St. Paulus eine Predigt, dadurch die Sünde offenbaret, die Herzen erschreckt und getödtet werden.

Um dieser Ursache willen hat Gott seinen Sohn gesandt und unter das Gesetz gethan, daß er uns von solcher Last durch sein Genugthun erlösete, und durchs Evangelium wieder tröstete und lebendig machte. Wie denn das Neue Testament seine eigene und neue Pfingsten hat, darin Gott vom Himmel sich anders hat hören lassen, denn hie; nicht mit Donner und Blitz, sondern mit einem feinen sanften Wind und mit schwebenden Feuerzungen; dafür erschraken die Herzen nicht, wie am Berge Sinai, sondern wurden fröhlich und muthig, und preiseteten und lobeten Gott, unangesehen, was die Welt dazu that.

\section{Votum}

GOtt der Vater erleuchte und regiere unsere Herzen durch den heiligen Geist, daß wir nach seinen heiligen Zehn Geboten heilig leben, und um Jesu Christi willen an Leib und Seel gesegnet und selig werden, Amen.

\chapter{Das 20. Kapitel}

\section{Vorrede}

Unter den größten Dingen, die im menschlichen Geschlecht geschehen sind, ist wahrhaftig eines, die Offenbarung Gottes, da er mit himmlischer Stimme sein Gesetz oder die heiligen Zehn Gebote wiederholet hat auf dem Berg Sinai. Für solche Offenbarung sollen wir allezeit, und also heute abermal Gott danken; sollen auch Gottes Wort für wahrhaftig halten und Gottes Geboten folgen, mit Thun und Lassen, in allem, was darinnen geboten und verboten ist. Denn Gottes Gesetz ist ein Kerker, geordnet zu schützen die auswendige Zucht und Ehrbarkeit; und die Zehn Gebote sind wie zehn Ketten und Fesseln dadurch unsere
\quellseite{152}
Glieder gebunden werden, daß sie sich nicht muthwillig an Gott und dem Nächsten vergreifen. Darnach, Gottes Gebote sollen auch unser Kompaß, Richtschnur und Winkelmaß sein, darnach wir ein christlich und unsträflich Leben führen sollen. Darauf höret davon an das 20. Kapitel.

\section{Mose empfängt die heiligen zehn Gebote GOttes}

a Das ist, was du zu thun hast.

\section{Summarium}

Ist eine schöne Predigt, welche Gott selbst vom Himmel herunter gethan hat, in welcher er seinem Volk anzeiget, wie sie sich halten sollen, wenn sie nach seinem Willen leben wollen. B. 1. Aber hie müssen wir unsere verderbte Natur kennen lernen. Denn weil wir solchen Willen Gottes wissen, und doch denselben weder wollen noch können thun: daher könnens, daß wir gegen
\quellseite{153}
Gott nimmermehr kein gut Gewissen können haben, wir müssen uns für Sünder bekennen und Gottes Zorn, Ungnade und Strafe gewarten. Und ist hie kein ander Mittel, dadurch uns zu helfen ist, denn daß wir Vergebung durch Christum glauben. Sonderlich aber ist hie zu merken, daß Gott so fleißig verbeut, man soll keinen Gottesdienst fürnehmen ohne und außer seinem Wort und Befehl.

\section{Votum}

Lob, Ehr, Preis und Dank sei GOtt für seine heiligen Zehen Gebote; der lehre uns durch seinen heiligen Geist dieselben nicht allein mit dem Munde zu bekennen, sondern daß wir auch nach denselben allzeit gehorsamlich mit der That leben, Amen.

\chapter{Das 21. Kapitel}

\section{Vorrede}

An Unterricht, wie man gegen Gott und den Menschen leben solle, mangelts nicht; daran mangelts, daß wenig solchem Unterricht gemäß leben. Derowegen laßt uns Gott danken, daß er uns lehret, wie wir ihm mit Thun und Lassen schuldigen Gehorsam leisten sollen; und schaue ein jedes zu, daß es also lebe, auf daß es nicht als ein muthwilliger Uebertreter Gottes Gebots sterbe und verdammt werde. Denn sind doch die unvernünftigen Kreaturen Gott ihrem Schöpfer gehorsam: als auf Gottes Befehl stehen die Wasser und Sonne stille, der Himmel gibt Himmelsbrod und Wachteln. Der harte dürre Steinfels gibt Wasser. Der Wallfisch speiet Jonam aus. Das Feuer muß die drei Männer im feurigen Ofen nicht verbrennen. Die Löwen müssen Daniel nicht fressen. Sollten wir vernünftigen Kreaturen nicht viel mehr Gott gehorsam sein? Müßten wir doch härter denn die Felsen, wilder denn die Löwen, ja gräßlicher, denn der Teufel selbst sein, der auf Christi Befehl weichet und ihm gehorsam leistet. Darauf höret nun an das 21. Capitel.

\section{Ordnung vom Gesinde, Item, über Mörder und Todtschläger}
\quellseite{154}

\section{Summarium}

Es ist eine Ordnung der Gerichtshändel, so in Todtschlägen oder in andern Verletzungen bei den Juden haben müssen gehalten werden. Solche Ordnungen gehen uns weiter nichts an, denn sie den natürlichen Rechten gemäß sind, was außer demselben zum Regiment verordnet ist, das mag man nach jedes Landes und Ortes Gelegenheit setzen oder fallen lassen. Sonderlich aber ist es fein, daß Gott verbeut, es soll keine Freiheit noch Geleit die helfen, so fürstlich und mit List einen Todtschlag begehen, V. 14. Item, daß die Kinder, so Vater und Mutter entweder schlagen oder ihnen fluchen, sollen getödtet werden, V. 15. Solche Satzungen stimmen mit den natürlichen Rechten; derhalben sie sollen gehalten und solche unartige Leute nicht geduldet werden.

\section{Votum}

Gott steure den unmorrischen Teufeln und blutgierigen Menschen, und verleihe uns, daß wir nach allen seinen Geboten thun und lassen, was er uns geboten und verboten hat, Amen.

\chapter{Das 22. Kapitel}

\section{Vorrede}

Wir sollen nicht allein Hörer, sondern auch Thäter sein des Wortes Gottes. Denn wir unserm lieben Gott in der heiligen Taufe Gehorsam gelobt und geschworen haben. So ist er auch unser zu Leib und Seel, Ehr und Gut mächtig, und wir ohne ihn nicht einen Finger regen, oder einen Odem holen, viel weniger größere Thaten üben können. So daß auch reichlich und wohl um uns Gott mit viel und großen Wohlthaten verdienet, daß wir ihm zur Dankbarkeit Gehorsam leisten. Darwegen helfe Gott, der heilige Geist, uns allen, daß wir also Gottes Wort anhören, daß wir nach seinen Geboten christlich leben auf Erden, damit wir nach seinen gnädigen Verheißungen leben in Seligkeit im Himmel. Darauf hört ferner an das 22. Kapitel.

\section{Strafe des Diebstahls und anderer Sünden}
\quellseite{155}

a Götter) heißen die Richter, darum, daß sie an Gottes Statt nach Gottes Gesetz und Wort, nicht nach Eigendünkel, richten und regieren mußten; wie Christus zeuget.

b Dringen) dringen und wuchern sind zweierlei: dringen ist, wenn du deinem Nächsten zwingest, zu bezahlen mit deinem Vortheil und seinem Nachtheil. Wuchern weiß man wohl, was es sei.

c Fülle) heißet er alle harte Früchte, als da sind Korn, Gerste, Aepfel, Birnen, da man Speise von macht.

d Thränen) heißet er alle weiche Früchte, da man Saft und Tran von macht, als da sind Weintrauben, Oele.
\quellseite{156}

\section{Summarium}

Es ist eine Ordnung der Gerichtshändel, so sich zutragen im siebenten Gebot, mit Diebstahl. Darnach sind viel seiner bürgerlicher Satzungen, daß man die Fremdlinge nicht übersehen, V. 22., Wittwen und Waisen handhaben, armen Leuten ohne Wucher leihen und ihr nöthiges Pfand oder Hausrath, daß sie zur Nahrung nicht gerathen können, über Nacht nicht behalten soll, V. 25. Denn mit solchem Leihen würde ihnen übel gedienet. Item, daß man der Obrigkeit nicht fluchen, noch sie lästern, sondern sie fürchten und ehren soll.

Item, daß man den Gottesdienst und die Personen, so dazu verordnet, mit dem, daß jedermann Gott bescheret hat, gern fördern und erhalten, die Gotteslästerer aber mit dem Tode strafen soll, V. 28.

\section{Votum}

Gott versiegele sein Wort in unsern Herzen, und helfe, daß wir uns gegen ihn kindlich, und gegen unsern Nächsten brüderlich allzeit verhalten, in Kraft des heiligen Geistes, Amen.

\chapter{Das 23. Kapitel}

\section{Vorrede}

Die zehen Gebote Gottes sollen aller Menschen täglicher Sündenspiegel und Laftafel sein, also, daß ein jeder seine Sünde daraus erkennen lerne, und die Gott abbitte, in wahrer Reue und Glauben. Und weil die Gebote Gottes lehren, was man dagegen Gutes thun soll, daß vom Bösen bewegen ein jedes lasse, und in guten Werke sich übe. Daher singen wir: Das Gesetz darum ist gegeben, daß es sei ein Spiegel zart, darinnen wir sehen unsre sündige Art, die in unserm Fleisch ist verborgen. Und: Die Gebote all uns gegeben sind, daß ein jedes soll lernen wohl, wie man für Gott leben soll. So verleihe Gott, daß wir alle seiner Gebote also täglich recht und christlich gebrauchen, Amen. Davon predigt auch das 23. Kapitel.

\section{Von Festen und Feiertagen}

a Versäume gern] Id est, relinquo mea interim et damnum accipio. Das ist: ich versäume indessen das Meine und leide Schaden.

b (Gedenken) ihr sollt von keinen Heiligen predigen, als von Göttern, die euch helfen, noch dafür danken, sondern Gott allein. Denn gedenken heißt so viel als predigen, rühmen, danken, Gottesdienst pflegen, ut: hoc tacite in mei commemorationem. 3. B.: daß thut zu meinem Gedächtnis.

c Dreimal) das ist das Osterfest im April, Pfingsten im Brachmonat, und das Laubhütten-Fest im Weinmonat. Davon lies im 23. Kapitel des dritten Buches Mose.
\quellseite{157}

d Des Jahres Ausgang heißt er den Weinmonat, daß alsdann aus ist mit Frucht wachsen und sammeln.

e Das Blut rc.) das ist, du sollst das Opferlamm nicht schlachten noch opfern, ehe denn alles gesäuerte Brod aus deinem Hause kommt.

\section{Summarium}

Ist zum Theil eine Auslegung des achten Gebotes, daß man wahrhaftig sei, recht richte und mit Geld und Gut sich nicht überwiegen lasse. Darnach ordnet er die drei fürnehmsten Feste im Jahr, B. 14., da alle Juden zum Tabernakel und Tempel kommen mußten; und sonderlich gebet Gott, man soll vor ihm nicht leer erscheinen, B. 16., das ist, man soll zur Kirche und zur Erhaltung des rechten Gottesdienstes gern geben, und nach Vermögen dazu helfen.

\section{Summarium}

Ist eine schöne Vermahnung, damit er die Predigt von den zehen Geboten beschließt. Auf daß, weil Gott durch sein Wort seinen Willen hat offenbaret, wir demselbigen sollen folgen. Thun wir es, so wolle er Brod und Wasser, das ist, Essen und Trinken segnen, alle Krankheit abwenden, gesundes Alter verleihen und allenthalben segnen, B. 25, 26. Thun wir es aber nicht, so soll die Strafe auch nicht außen bleiben.

Da stehet man nun, was es für Ursache sei, daß so viel arme, kranke, unglückselige Leute auf Erden sind: nämlich, daß der meist Theil solches durch eigene Sünde und gottlos Leben verdienet. Darum, wer solchem Unglück entlaufen, und wie der 34. Psalm B. 13. sagt, gute Tage haben will, der fürchte den HErrn, thue Gutes und fliehe das Böse, so will Gott wiederum thun, was er soll und wir begehren.

\section{Votum}

Der treue, barmherzige GOtt helfe, daß wir ja seinem Willen nicht widerstreben, sondern nach demselben allzeit in Gehorsam christlich leben und selig sterben, Amen.
\quellseite{158}

\chapter{Das 24. Capitel}

\section{Vorrede}

Gleich wie bei dem Berge Sinai stund sechsmal hundert tausend Mann und hörten die Gesetzpredigt Gottes, zitterten und bebten, gelobten auch an, sie wollten Gott gern gehorsam sein in allem; aber es ward darnach bald vergessen, und geriethen in allerlei Sünden, daß nur zween kamen in das gelobte Land: also gehet es stets. Nicht alle, die getauft sind, nicht alle, die Gottes Wort hören, nicht alle, die da gebrauchen das Sacrament des wahren Leibes und Blutes des HErrn Christi, auch nicht alle, die sich des heiligen Geistes rühmen, erlangen das ewige Leben. Denn viele ihr Fleisch und Blut sich auch verführen lassen, leben und sterben in Sünden. Derowegen sollen wir zusehen, wie wir fürsichtiglich wandeln, und sollen verständig sein, was da sei des HErrn Wille, und sollen zugesagten Gehorsam Gott, nach seinem Wort, treulich leisten. Davon höret auch an das 24. Capitel.

\section[Moses steiget, nach Bestätigung des Bundes, wieder auf]{Moses steiget, nach Bestätigung des Bundes, wieder auf den Berg Sinai}

a Seine Hand) er schreckte sie nicht mit Donner und Blitz, wie zuvor das Volk erschrecket ward, Kap. 20, 18.

\section{Summarium}

Ist eine Historia, wie das Volk sich verpflicht-, Gottes Gesetze und Willen zu halten, V. 4. Aber wie solches erfüllet sei, werden wir in der Historia weiter hören; daß auch St. Petrus zeuget in den Geschichten der Apostel Kap. 15, 10.: Das Gesetz sei ein Joch, welches weder die Apostel noch die Väter gehalten haben. Darum spricht die Epistel an die Ebräer: Der Bund, so durch der Farren Blut gemacht, sei nicht kräftig gewesen. Und hat eben solcher Ursache halben Gott einen andern Bund aufgerichtet, der nicht auf äußerlichem Thun und Lassen, sondern auf Gottes Gnade beruhen soll, durch das Blut unsers Seligmachers Jesu Christi, wenn wir an ihn glauben und seines Verdienstes uns trösten.
\quellseite{159}

\section{Votum}

Gott der heilige Geist erleuchte und regiere uns, daß wir auch thun, was wir Gott zugesagt haben, und daß wir alle Sünde hassen und Tugenden üben, auf daß wir zeitlicher und ewiger Strafe entfliehen, Amen.

\chapter{Das 25. Kapitel}

\section[Inhalt der folgenden sieben Kapitel: Von dem Heiligtum]{Inhalt der folgenden sieben Kapitel: Von dem Heiligtum und was mehr dazu gehöret hat}

\section{Vorrede}

Diese folgenden sieben Kapitel lehren anders nichts, denn wie Gott Mose auf dem Berge habe Befehl gegeben, von sonderlichem Gottesdienst aufzurichten, der unter den Juden bis auf Christus Reich sollte gehalten werden.

Denn hie ordnet er erstlich, wie man die Hütte oder Kirche bauen, V. 8., die Lade des Bundes, V. 10., samt dem Gnadenstuhl, V. 17., und Cherubim zurichten, V. 18., den güldenen Tisch mit den Schaubroden, V. 23., den güldenen Leuchter, V. 31., den Vorhang samt anderem Geräthe, das man dazu bedurfte, machen soll.

Darnach ordnet er die Priester, wie sie gekleidet sein, was und wenn sie opfern sollen, auch wie man sie weihen soll. Solches alles, ob es jetzt schon nicht mehr im Brauch ist, ist doch derhalben sein zu wissen, daß man unsern Papisten desto baß begegnen kann, welche auch ihre Kirchen, Pfaffen, Kelch, Meßgewand und andere Kirchengefäße weihen; aber doch mit einer weit andern Meinung, denn Moses.

Moses hatte Befehl von Gott, er sollte es so machen und anders nicht; darum weiset ihm auch Gott auf dem Berg ein Muster, da er alles vor ihm stehen sahe, wie es sein sollte.

Unsere Papisten aber haben außer und ohne Gottes Befehl Pfaffenweihe, Kirchenweihe und andere Gerätheweihe aus ihrem eigenen Kopf erfunden und fürgenommen, können auch keine andere Ursache anzeigen des Papstthums halben, daß der Papst soll das oberste Haupt der Kirchen sein; item der Pfaffen und Kirchenweihe etc., denn daß im Alten Testament also gewesen sei. Wo es nun im Neuen Testament nicht auch sollte sein, sagen sie, das Alte Testament wäre herrlicher, denn das Neue. Aber wir wissen eine andere Herrlichkeit, nämlich, daß der Sohn Gottes oberster Priester ist, und mit seinem Opfer uns versöhnet hat, der weihe seine Kirche mit dem heiligen Geist und gibt tüchtige Kirchendiener. Derohalben soll solches päpstliches Wesen durchaus von niemand anders gedacht werden, denn wie andere Menschenkünste und Ordnung, so ohne alles Gewissen mögen, ja sollen geändert werden; sonderlich, wenn man sie für nöthig achten will, wie der Papst thut. Dies ist das erste, das man bei diesen Kapiteln soll merken.

Das andere ist, daß Gott allen solchen Unkost allgemein auf das Volk schlägt, daß sie zusammen schließen und ein jeder das Seine geben soll, daß es also möge gebauet und angerichtet werden, V. 20. Uns zur Lehre, daß wir auch gerne dazu geben und helfen sollen, was zur Förderung des rechten Gottesdienstes und Kirchen-Ämtern vonnöthen ist.

Das dritte ist die heimliche Deutung mit solchem Bau und Ordnung. Die Hütte ist die christliche Kirche, in derselben liegt das Zeugniss, das ist, die heilige Schrift und das Wort Gottes.

Der Gnadenstuhl darin, das ist unser lieber Herr Christus Jesus, durch welches Verdienst wir mit Gott versöhnet sind. Derselbe ist auch der Hohepriester, welcher sich geopfert hat für unsere Sünde, und ist von seinem himmlischen Vater gesalbet oder gewehet mit seinem heiligen Geist, welchen er ausgegusst durchs Evangelium, daß wir auch gesalbte Priester werden und Gott dienen durch den rechten Glauben, mit der Predigt von seiner Güte und Gnade, mit reinem Gebet, mit Opfern unser selbst und mit täglichem Töden des alten Adams.

Solches alles sind die sonderlichen herrlichen Gaben des Neuen Testaments, welches Gott den Juden durch solchen Brauch und Ordnung, gleich als durch ein Fürbild hat anzeigen wollen; wie die Epistel an die Ebräer auch lehret. Dieses sei zu einer einfältigen Anleitung von diesen Kapiteln genug gesagt.

\section{Freiwillige Opfer zur Stiftshütte}

a Fürbild) das ist der Hauptstrich einer, darin Mose zeuget, daß sein Regiment soll aufhören und nicht das rechte ewige Wesen, sondern ein Fürbild und Beispiel sein des Reiches Christi. Und ohne allen Zweifel hat hieran Christus mit den zwei Jüngern zu Emmaus geredet, Luc. 24, 27.
\quellseite{160}

b Dir zeugen) das ist, dabei, als bei einem gewissen Zeichen und Zeugniß, will ich dich wissen lassen, daß ich da bin gegenwärtig, daß ich daselbst reden werde.

c Centner) dreißig Pfund Goldes.

d Und siehe zu) er wiederholts hie zum andern Mal, daß er droben in diesem Capitel sagt: Es soll alles viel Fürbild und Figur sein mit Mose, auf daß es, zweimal gesagt, ja gewißlich geglaubt werde.

\chapter{Das 26. Kapitel}

\section{Bau der Stiftshütte}
% Omitted: numbered Bible verse, excluded; printed folio or catchword marker; running head; not text apparatus.
\quellseite{162}

\chapter{Das 27. Kapitel}

\section{Vom Brandopferaltar, Vorhof und heiligen Oel}

a Hörner) aufgerichtete kleine Säulen, oben mit Knöpfen oder Blumen.

b Moed) das hebräische Wort Moed haben wir nicht anders wissen noch wollen deuten. Es soll aber so viel heißen als ein gewisser Ort oder Stätte, wie eine Pfarrkirche oder Stift, dahin das Volk Israel kommen und Gottes Wort hören sollten, damit sie nicht ihrer eigenen Andacht nach hin und wieder liefen, auf Bergen, in Gründen und andern Oertern, Gott zu opfern.

\chapter{Das 28. Kapitel}

\section{Priesterliche Kleiderzierde Aarons und seiner Söhne}
\quellseite{163}

a Licht) was das gewesen sey, weiß man jetzt nicht mehr. Hebräisch heißt es Urim und Thummim, 4 Mos. 27, 21.
\quellseite{164}

b Füllen) dies Füllen ist eine hebräische Sprache, der man muß gewöhnen. Und war das, wie im folgenden Capitel stehet, daß in der Weihe den Priestern die Hände mit Opfer gefüllet wurden für dem HErrn.

\chapter{Das 29. Kapitel}

\section{Einweihung und Opfer der Priester}
\quellseite{165}

a (Ein Feuer des HErrn) das Wörtlein „Feuer des HErrn“ brauchet Mose sehr viel, und wir sollten's gewöhnen, daß es so viel heiße, als ein Opfer des HErrn. Als wenn du von ferne ein Feuer sähest, und fragest, was ist das? So spräche man, es ist Gottes Feuer, das ist, man opfert dasselb Gott. Per synecdochen vel aliam figuram.

b Entsündigen) das ist, absolviren und lossprechen, wie Ps. 51, 9. Asperges me Isopo, das ist: Entsündige und absolviere mich mit Isop.

c Hin) ist fast eine halbe Metze.
\quellseite{166}

\chapter{Das 30. Capitel}

\section[Vom Räuchaltar, der aufgesetzten Steuer, Handfaß, Salböl und]{Vom Räuchaltar, der aufgesetzten Steuer, Handfaß, Salböl und Räuchwerk}

a Dieser Altar ist außen für dem Fürhang im Heiligen, und nicht hinter dem Fürhang im Allerheiligsten gestanden.

b Lampen zurichtet) Benefacit, depurando Evangelium ab opinionibus, daß es sein brennen kann. Das ist, wenn man das Evangelium von Menschenfälschungen reinigte.

c (Gera) gilt bei uns drei Löwenpfennige, oder ein Dreigröschlein. Ein halber Seckel ist ein Ort eines Gülden.
\quellseite{167}

\chapter{Das 31. Kapitel}

\section{Bestellung der Werkmeister. Sabbathfeier. Gesetztafeln}
\quellseite{168}

\chapter{Das 32. Kapitel}

\section{Vorrede}

Gott belohnet das Gute und strafet das Böse; und auf große Sünden folgen große Strafen, und solches ist die Belohnung der schändlichen Sicherheit. Das lehret Gott, da er Ezech. 8. also sagt: Ich will wider sie mit Grimm handeln, mein Auge soll ihrer nicht verschonen und will nicht gnädig sein. Und abermal Malachi am 3. Kapitel also: Ich will ein schneller Zeuge sein wider die, so mich nicht fürchten, denn ich bin der HErr, der nicht leuget, spricht der HErr Zebaoth. So ist dies auch wahr: Wie ein Bock härter stößt, wenn er weit aushole[t]: also strafet Gott desto schwerer und heftiger, wenn er die Strafe aufschiebt. Derowegen, weil ohne das Jammer und Noth genug ist in allen Winkeln der Welt, in allen Ständen, so laßt uns Sicherheit fliehen, beten und gottesfürchtig leben, auf daß wir mit unsern Sünden Gott nicht mehr erzürnen und zu Strafen weiter Ursache geben. Darauf höret auch an das 32. Kapitel.

\section{Abgötterei mit dem gegossenen Kalbe}

a Verzog das redet Moses mit einem Wort also, als habe er das Volk verlassen mit Schanden, und sey gestorben, daß sie nicht wissen, wie sie nun thun sollen. Gleich als wenn uns Gott eine kleine Zeit verläßt, denken wir, er lasse uns in Schanden stecken, müssen anderwärts Hülfe suchen; da wird denn solch Kalb unser Gott.

b Entwarf das ist, er machte es ihnen für, was sie für ein Bild machen sollten. Das bedeutet, daß Menschenlehre dem Volke vorbilden, was sie für Werke thun sollen, damit sie Gott dienen. Denn hie siehst du, daß sie in diesem Kalbe vermahnet haben dem rechten Gott zu dienen, weil Aaron rufen läßt: Es sei des HErrn Fest, und bauet ihm einen Altar.

\section{Summarium}

Dies ist eine erschreckliche Historia, wie das Volk Israel, welches so große Wunder in Egypten und in der Wüste gesehen und erfahren, und sich gegen Gott, seine Gebote zu halten, so theuer verpflichtet hatte, so bald alles vergessen und in greuliche Abgötterei gefallen ist und ein gegossen Kalb angebetet hat, V. 4.
\quellseite{169}
Uns zum Exempel, wie es Paulus selbst deutet, 1 Cor. 10, 6., daß wir nicht sicher noch vergessen sein, sondern in Gottes Furcht stehen sollen und lernen, daß es, so bald Gott die Hand abzieht, mit uns aus ist. Darum ist vonnöthen, daß wir stets bitten, daß Gott durch seinen Geist und Wort uns erhalten, und nicht in Anfechtung führen wolle.

Darnach ist ein fein Exempel an Mose, wie er für das Volk bittet und das Gebet allein auf Gottes Güte und Verheißung stellet, und Gott solches annimmt, V. 11.

c Das Wort Phrea lautet und heißt schier dem Deutschen gleich, frei; will so viel sagen: Aaron hatte das Volk frei, los und ledig gemacht von Gottes Wort und Gehorsam, wie die Menschen-Lehre thut, machet das Volk wild, frei, los und bereit zu aller Abgötterei, und thut doch der Meinung, als wollten sie den Leuten aufhelfen, und besser machen. Und ist doch (wie es im Hebräischen heißt) ein Schemez à Schama, ein Gehör oder Geschwätz, da nichts hinter ist.

d Nemo evadet poenam peccati sui. Das ist, Niemand wird der Strafe seiner Sünde entgehen.

\section{Summarium}

Ist eine Historia, wie Moses um des Aergernisses willen heißt die Kinder Levi erwürgen und erstechen, was sie antreffen, V. 27. Auf daß er dem Zorn und der Strafe Gottes fürkomme, und den größten Haufen errette. Denn solches gehöret der Obrigkeit zu, daß sie öffentliche Aergerniß strafe; auf daß Gott nicht zu größerem Zorn beweget werde und den ganzen
\quellseite{170}
hausen darum angreife, so man öffentliche Aergernisse ungestraft läßt. Das aber Moses die steinernen Tafeln zerbricht, V. 19. ist eine Deutung, daß menschliche Natur das Gesetz Gottes nicht vermag zu halten, es komme denn Gott, und helfe nicht allein mit seinem heiligen Geiste, sondern auch mit Vergebung der Sünden, um Christi willen.

\section{Votum}

GOtt sei uns allen gnädig und barmherzig, und helfe, daß wir für Abgötterei und allen andern Sünden uns hüten, und nach seinen heiligen Geboten gehorsamlich bis zu unserm seligen Todesstündlein leben, Amen.

\chapter{Das 33. Kapitel}

\section{Vorrede}

Kein recht Gebet ist verloren. Denn Gott erhöret alle unsere Seufzer und schreibet sie auf; und wenn wir nur anfangen zu schreien: Vater, so gibts in dem Himmel einen mächtigen starken Schall: Sieh bin ich mein Sohn, meine Tochter, was willst du? Derhalben ist ein christlich Gebet gar sehr kräftig, mächtig und geschäftig, und es dringet durch alle Himmel, und richtet aus drüber und nicht drunter, denn wir je gebeten haben. Daher das gläubige Gebet wohl mag genennet werden der edelste Schatz der Christenheit, die beste Wehre und der höchste Trost, und die bewährteste Arznei in Nöthen, und der schönste Gottesdienst. So laßt uns alle Tage oft und fleißig im Geist und in der Wahrheit beten, Gott anrufen und danken. Dies erinnert auch das 33. Kapitel.

\section[Mose bittet für das Volk, und begehret des]{Mose bittet für das Volk, und begehret des Herrn Herrlichkeit zu sehen}

a Den Rücken Mose sehen alle Werckheilige, die das Gesetz nicht verstehen, noch unter Augen kommen.

b Angesicht, heißt die Wolke und feurige Säule, darin Gott gegenwärtig war.
\quellseite{171}

c Mein Angesicht) also heißt Gottes Angesicht nicht die Wolke noch Säule, sondern er selbst, wie er spricht: Der Mensch wird nicht leben, so mich siehet. Aber es ist alles von Christo gesagt, welcher sollte erscheinen in der Menschheit hernach, wenn Moses Reich ein Ende hätte. Indes soll Moses im Fels stehen und den vorbeygehenden Christum sehen und predigen, bis er komme; daß also Israel Gottes Wort immer habe, bis auf Christum.

\section{Summarium}

Ist eine feine Historia, wie Moses nirgend hinziehen will, es ziehe denn Gott mit, B. 15. 16. Solches soll eine christliche Obrigkeit allzeit thun, daß sie durch's Gebet Gott zu einem Vorgeher, Geleitsmann und Gehülfen nehme; sonst wird wenig Glück dabei sein, wo Gott mit uns nicht dabei ist.

Daß aber Moses begehren, Gottes Herrlichkeit zu sehen, und Gott sich solches weigert, B. 18. 19., ist uns zur Lehre fürgeschrieben, daß wir uns an der Offenbarung, so Gott in Christo und durch sein Wort uns hat widerfahren lassen, sollen genügen lassen. Denn außer diesem Anblick der Gnaden, welchen wir in Christo und im Wort haben, kann man Gottes Herrlichkeit nicht sehen und leben. Unsere Natur kanns nicht erdulden noch tragen, bis in jenem Leben, da das Fleisch nicht mehr wird sterblich sein, da werden wir Gott von Angesicht zu Angesicht sehen.

\section{Votum}

Der treue, barmherzige Gott sei und stehe auch bei uns in diesen gefährlichen Zeiten, blicke uns an mit Gnaden unter unserm Kreuz, und helfe, daß wir seine Herrlichkeit sehen in alle Ewigkeit, Amen.

\chapter{Das 34. Kapitel}

\section{Vorrede}

Ein jeder Mensch soll täglich gedenken an Gottes Gerechtigkeit und Barmherzigkeit. An Gottes Gerechtigkeit sollen wir gedenken, auf daß wir uns für Sünden hüten; an seine Barmherzigkeit aber, auf daß wir wider die Sünde Trost haben. Denn wie Gott, als ein eifriger Gott, die Sünde an allen Unbußfertigen strafet: also vergibt er die Sünde, als ein barmherziger Gott, allen, die ihn darum bitten. Von diesen beiden lehret also die heilige Schrift: Jeremias sagt im 5. Kapitel: Eure Mißthat und Sünden wenden das Gute von euch. Und David im 103. Psalm spricht: Wie sich ein Vater über Kinder erbarmet, so erbarmet sich der HErr über die, so ihn fürchten. Darauf höret weiter an das 34. Kapitel.

\section{Neue Gesetztafeln. Bund Gottes. Moses glänzendes Angesicht}

a Predigt) Mose berichtet hie das Geheimniß, daß Christus der HErr sey, der für ihn (das ist, für den Volk des Geistes) werde übergeben und predigen, daß alle Welt sündig, allein aus Gnaden selig werden müsse, wie er lehret Kap. 33, 19. sagt: Wem ich gnädig bin, dem bin ich gnädig, das ist, es soll mirs keiner mit Werken abverdienen, es soll Gnade sein, und kein ander. Wie denn des Evangelii Predigt gleich also lehret.
\quellseite{172}

b Halsstarrig Volk) Opera nostra non faciunt, si misericordia non comitetur. Das ist: unsere Werke werden nicht thun, wenn uns nicht Barmherzigkeit begleitet. Wir sind wohl so böse.

\section{Summarium}

Ist eine Predigt, in welcher die zehn Gebote wiederholet, sammt etlichen andern Geboten von den fürnehmsten Festen. Und ist hie sonderlich zu merken, wie das erste Gebot wird heraus gestrichen, daß Gott barmherzig, gnädig, geduldig, und von großer Gnade und Treue sei, die Sünde vergeben wolle, weil wir doch für ihn nicht können unschuldig sein. Darum soll man sich solches Gottes allein trösten und auf ihn trauen. Wiederum ist er ein eifriger Gott, der Abgöt%
terei,\quellseite{173} Mißglauben und gottlos Wesen nicht leiden, sondern bis ins dritte und vierte Glied strafen will, V. 7. Darum soll man solchen Gott allein fürchten und sein Wort und Willen halten.

\section{Summarium}

Ist eine Historia, wie Moses, da er vom Berge gegangen und das Gesetz gebracht hat, unter seinen Augen geleuchtet habe, daß ihn niemand hat ansehen können. Wie denn solcher Ursache halben die Maler ihn mit Hörnern malen. Es sind aber nicht Hörner, sondern ein Glanz gewesen, der von seinem Angesicht gegangen und den Menschen ist unleidlich gewesen, derhalben er eine Decke hat müssen fürhängen.

Solches ist eine seine Deutung, wie es St. Paulus 2 Cor. 3. angzeigt, daß man das Gesetz nicht recht verstehen noch thun kann, wo der Geist Gottes nicht recht im Herzen ist, welcher allein durch Christum uns geschenkt wird.

Denn gleich wie die Juden nur den Fürhang für Moses Angesicht haben gesehen und den Glanz nicht leiden noch sehen können: also sind alle die, so sich unterstehen das Gesetz zu thun aus eigenem Vermögen. Die Werke mag man wohl nachthun dem äußerlichen Wesen nach; aber das Herz bleibt unrein, einen Weg wie den andern. Wo aber Gottes Geist durchs Evangelium ist, da werden die Herzen geändert und das Gesetz recht erfüllet. Denn ob schon dem Werken etwas fehlet, wird es doch durch Vergebung der Sünde in Christo erstattet.

\section{Votum}

GOtt gebe und erhalte in uns allen seinen heiligen Geist, daß wir durch ihn die heiligen zehn Gebote recht verstehen lernen, und daß wir nach denselben allzeit in Heiligkeit und Gerechtigkeit leben, Amen.

\chapter{Das 35. Kapitel}

\section{Vorrede}

Wie kein Mensch nimmer aus der Kirche gehen soll, er habe denn zuvor gebetet, und etwas Nützliches aus Gottes Wort gelernt: also welchen Gott gesegnet hat, wenn er aus der Kirche geht, der soll sein Almosen im Kasten Gottes auch einlegen, den Hospitalen und hausarmen Leuten zur Beisteuer. Wohl denen, die solches thun aus wahrem Glauben. Denn Gott wird an ihrem Leib, Seele und Nahrung zeitlich bleiben und im künftigen Leben aus Gnaden mit ewiger Glorie, Freude und Herrlichkeit vergelten. Davon lehret der heilige Prophet Jesaias im 58. Kapitel, da er spricht: Entzeuch dich nicht von deinem Fleisch (das ist, verlasse deinen Nächsten nicht mit Almosen und anderer nöthiger Hülfe): alsdann wird dein Licht hervor brechen, wie die Morgenröthe, und deine Besserung wird schnell wachsen, und deine Gerechtigkeit wird vor dir hergehen, und die Herrlichkeit des HErrn wird dich zu sich nehmen. Dies lehret auch das 35. Kapitel.

\section{Sabbathsruhe. Freiwillige Steuer. Beruf der Werkmeister}

a Was Gott droben vom 25. Kapitel an bis aufs 32. Mose befohlen hat von der Wohnung und alle ihrem Geräthe zu machen, das zeiget Mose dem Volk an in diesem Kapitel.
\quellseite{174}

\section{Summarium}

Dies ist ein schön Exempel, in welchem wir sehen, wie die Juden so reichlich und gern Gold und Silber, und was sie haben, dazu geben, daß die Hütte oder Kirche gebaut und der rechte Gottesdienst gefördert und angerichtet würde, V. 21. Und zur Lehre, daß wir mit dem, so uns Gott gegeben hat, auch gern sollen helfen, auf daß das Predigtamt, die Schulen, sammt anderem, so zur Unterhaltung des Kirchendienstes vonnöthen, möge erhalten werden. Denn wer solches thut, der bauet recht an Gottes Tempel und leget sein Geld also an, daß Gottes Ehre und der Menschen Seligkeit dadurch gefördert werde.

Vom 10. Vers bis auf den 20. werden nur
\quellseite{175}
diese Worte gelesen: Und wer unter euch verständig ist, der komme, und mache, was der HErr geboten hat.

\section{Votum}

Helffe der treue barmherzige Gott, daß all unser Thun und Lassen zu seinem Ehren, und unsers Nächsten Nutz und Wohlfahrt gelange, und daß ein jedes mit Gebet und Hülfe, zur Erhaltung der wahren Gottesdienste, treulich diene, Amen.

\chapter{Das 36. Kapitel}

\section{Vorrede}

Diese vier Kapitel ist nicht vorzunehmen, daß mans von Wort zu Wort lese. Denn es ist nur eine Erzählung, wie man die Kirche sammt aller Zugehörung gebauet habe. Da können wir keine andere Lehre aus nehmen, denn diese: daß alles, was zur Kirchen und Gottesdienst im Alten Testament gebräucht hat, allein nach Gottes Wort, wie es Gott befohlen hat, ist gemacht worden. Denn eigen erwählter Gottesdienst außer und ohne Gottes Wort kann und will Gott nicht leiden. Darum gehet alle Klage in den Propheten darüber, daß das Volk auf Gottes Befehl nicht thun, sondern ein jeder ihm selbst ordnen und fürnehmen will, wie er Gott dienen möge.

Gleich wie unter dem Papstthum auch geschehen ist, da sind die rechten Gottesdienste, Glaube, Liebe, Beten, im Kreuz und Leiden geduldig sein, dahinten blieben, deren hat niemand gedacht. Aber Mönche und Nonnen werden, sonderlich kleiden, fasten, zu heiligen Wallfahrten gehen und dergleichen, davon Gott nicht ein Wort befohlen hat, dieß hat man für sonderliche Gottesdienste gehalten und fürgenommen; so doch ein jeder Hausvater in seinem Haus das für keinen Dienst achtet, noch ihm gefallen läßt, was sein Knecht ohne Befehl aus eigenem Willen fürnimmt; sondern das heißt dem Hausherrn gedienet, so man thut, was er befohlen hat.

Darum beschreibet hie Moses nach der Länge wiederum, auf daß seine Juden und wir lernen, wie er alles, was zu machen gewesen ist, nach Gottes Wort und Befehl hat machen lassen; auf daß wir in allen Dingen auch auf Gottes Wort sehen, und demselben Wort nach alles handeln und anrichten. Das ist die fürnehmste Lehre dieser vier Kapitel.

Darnach ist zu merken, wie Moses hie anzeiget, daß es eine Gabe Gottes sei, wenn einer ein Handwerk wohl kann, V. 1. Darum soll man es erkennen, und Gott dafür danken, und solcher Geschicklichkeit nicht mißbrauchen, wie doch gemeiniglich geschiehet. Denn wo sonderbare Gaben sind, da denket niemand, daß er zu Gottes Reich und des Nächsten Noth mit helfen soll; sein selbst Ehre, Geld, Gut und anderes suchet schier jedermann damit. Es versündigen aber die Leute sich also dadurch, daß Gott weder Glück noch Segen dazu geben will. Wie man stehet an vielen Künstlern, daß sie gemeiniglich arme Bettler sind, nicht der Kunst halben, sondern daß sie ihr nicht recht brauchen.

Zum dritten wird auch hie wiederholet, wie das Volk so gar reichlich habe zugetragen, mit Gold, Silber und anderem, daß Moses verbieten muß und heißen aufhören, V. 6. Uns zum Exempel, daß wir auch gern und reichlich zum Gottesdienst geben, und zu Unterhaltung der Kirchen willig helfen sollen.

\section{Die Stücke des Heiligthums werden verfertiget}
\quellseite{176}

\chapter{Das 37. Kapitel}

\section{Von etlichem Geräthe des Heiligthums}
\quellseite{177}

\chapter{Das 38. Kapitel}

\section{Ansehnliche Summe des angewandten Goldes, Silbers und Erzes}

a Weibern) diese Weiber waren die andächtigsten Wittwen und Weiber, die mit Fasten und Beten für der Hütte Gott dieneten, wie 1 Kön. 2, 22. gezeuget. Und Paulus 1 Tim. 5, 5. beschreibt. Wie auch St. Lucas die heilige Prophetin Hanna rühmet, Luc. 2, 37.
\quellseite{178}

\chapter{Das 39. Kapitel}

\section{Priesterliche Kleiderzierde}
% Omitted: numbered Bible verse, excluded; printed folio or catchword marker; running head; not text apparatus.
\quellseite{180}

\chapter{Das 40. Kapitel}

\section{Stiftshütte aufgerichtet und eingeweihet}
\quellseite{181}

\section{Summarium}

Es ist eine Historia, wie Moses alles mit der Kirchen, und was dazu gehöret, angerichtet habe. Und ist solches, wie oben auch gemeldet, sein zu wissen, um des Anspiels willen, so der Papst und Bischöfe ohne Gottes Befehl mit Kirchen, Kleidern und Priesterweihen, auch mit Glockentaufen und andern angerichtet haben. Denn hie wiederholets der Text achtmal, da er allwege spricht: Moses habe solches gemacht, wie ihm der HErr geboten hat, B. 16. 19. 21. 23. 26. 27. 29. 32.

Weil nun solches alles eine mosaische Ordnung, und im Neuen Testament aufgehoben und abgethan ist, kann jedermann leichtlich abnehmen, daß solche \unsicher{Papstweihe}, ohne Gottes Wort fürgenommen, gleich so wenig zu halten oder zu achten sei, als andere Menschen-Fündlein.

\part{Das dritte Buch Mose}

\chapter{Das 1. Kapitel}

\section{Vorrede}

Das 1., 2., 3., 4., 5., 6., 7., 8. u. 9. Kapitel werden in der Kirche nicht gelesen.

Dies Buch hat den Namen in der lateinischen Kirche, und heißt Leviticus, darum, daß es fürnehmlich das Priestertum fället, und ordnet mit Gesetzen und Rechten, nach welchen die Priester mit dem Opfer sich halten und das Volk lehren sollen.

Da stehet man fürnehmlich, wie das Priestertum nur um der Sünden willen eingesetzt ist, daß dieselbigen dem Volk erstlich angezeigt und darnach wieder versühnet sollen werden. Also, daß eines Priesters Amt anders nichts ist, denn mit Sünden und Sündern für Gott umgehen. Wie denn unser lieber HErr Christus darum auch den Namen hat, Ps. 110, 4., und ein Priester heißt, daß er sich der Sünder angenommen und für die Sünde sich selbst geopfert hat am Stamme des Kreuzes.

Daß man aber einen rechten Verstand habe von dem Priestertum und dem Opfer des Alten Testaments, muß man zum Eingang wissen, daß das Priesteramt und ganze Gesetz, so fern es das äußerliche Regiment und Kirchenordnung belanget, dazu gedienet hat, daß es nicht allein einen äußerlichen Unterschied zwischen Heiden und Juden machte, sondern daß es auch eine Uebung wäre des äußerlichen Gehorsams gegen Gott, welcher Gehorsam keine Vergebung und zeitliche Belohnung hätte. Aber zur Gerechtigkeit, so für Gott gilt, und zu dem ewigen Leben hat solches Priesteramt und Opfer nicht gedienet. Nach dem Spruch Joh. 1, 17.: Das Gesetz ist durch Mosen gegeben, die Gnade aber und Wahrheit ist durch Christum worden. Darum hat es auch ein Ende nehmen und das Gesetz sammt den Priestern und dem Opfer aufhören müssen, da Christus kommen ist; weil es nur ein Schatten und Fürbild der künftigen Güter in Christo gewesen ist, wie zum Ebräer am 10. B. 3. stehet.

Denn der rechte Priester, unser HErr Christus selbst, der opfert für der Welt Sünde nicht Kälber noch Ochsen, sondern sein eigen Leib und Blut. Durch dies einige Opfer werden wir von Sünden erlöset, wenn wirs mit Glauben fassen und unser Vertrauen darauf setzen. Durch dies einige Opfer werden wir mit Gott wahrhaftig versühnet, daß wir zur ewigen Kindschaft wieder kommen, welche wir durch die Sünde verloren haben.

Daß nun im Gesetz auch Sündopfer gewesen sind, ist nicht dahin zu rechnen, als hätten sie die Sünden für Gott dadurch abgelegt und versühnet. Denn dazu gehöret, wie gesagt, nicht Kühe- und Kälber-Blut, sondern das Blut unsers HErrn Jesu Christi. Und haben die heiligen Propheten und Patriarchen alle, von Adam bis auf Johannes, in solchem Fall allein auf den verheißenen Weibes-Samen und sein Opfer gewiesen, daß man desselbigen sich annehmen und wider die Sünde und Verdammniß trösten soll. Daß aber außer diesem Opfer die Juden, als ein sonder Volk, ihre eigene Opfer gehabt, die man auch Sündopfer nennet, solche Opfer haben dazu gedienet, daß die Juden am ersten damit für der Gemeine bekennet haben ihre Uebertretung am Gesetze, und sich darnach durch solch Sündopfer gegen der Gemeine wieder versühnet. Denn wo sie solche Ordnung nicht gehalten hät%
ten,\quellseite{182} wären sie Heiden und nicht mehr Juden geachtet worden. Neben demselbigen haben die Gläubigen solche Sündopfer angesehen als ein Fürbild des rechten Opfers, dadurch sie und die ganze Welt, der Verheißung nach, von Tod und Sünden durch Christum sollten erlöset werden. Durch solchen Glauben haben alle Heiligen im Alten Testament, vor und unter dem Gesetz, Vergebung der Sünden erlangt und sind selig worden. Das sei genug vom Unterricht und Eingang dieses Buches, von dem Opfer und Priester-Amt, davon die ersten neun Kapitel melden.

\section{Gesetz der Brandopfer}

\chapter{Das 2. Kapitel}

\section{Gesetz der Speisopfer}
\quellseite{183}

\chapter{Das 3. Kapitel}

\section{Gesetz von Dankopfern}

a Speise) das vom Feuer aufgefressen wird.

\chapter{Das 4. Kapitel}

\section{Gesetz von Sündopfern}
\quellseite{184}

a Ärgerte) das ist, mit Lehren oder Leben ungefähr zur Sünde und Schuld Ursache gäbe.
\quellseite{185}

\chapter{Das 5. Kapitel}

\section{Gesetz vom Schuldopfer}
\quellseite{186}

a Geweiheten) Sacrilegium, raptum scilicet ex hostiis per impudentiam. Das ist, Vergreifung am Geweiheten, das nämlich entwendet ist von den Opfern aus Versehen.

\section{Kap. 6. v. 1}

\chapter{Das 6. Kapitel}

\section{Vom Brand-, Speis- und Trankopfer}

a Nec extinguantur, nec alienus sit.
\quellseite{187}

b Aber alle das) hie siehst du, daß Moses klärlich zweierlei Sünd-Opfer oder zweierlei Brauch des Sünd-Opfers setzet: eines, da man das Blut nicht in das Heilige bringet zu versühnen, solches möchten sie essen. Das andere, da man das Blut in das Heilige bringet zu versühnen, solches müßte man nicht essen, sondern außer dem Lager alles verbrennen. Davon in der Epistel an die Ebräer 13, B. 11. Quorum animalium sanguis. Et supra, cap. 4. v. 5.

\chapter{Das 7. Kapitel}

\section{Vom Schuld- und Dankopfer}

a Lob-Opfer, Dank-Opfer) diese zwei Opfer sind in ein Opfer gerechnet. Dank-Opfer heißt, wenn sie Schafe, Ochsen etc. geschlachtet haben. Lob-Opfer, wenn sie Fladen und Kuchen (wie ein Speis-Opfer) dazu gethan hatten. Und nennet also eins das andere, daß es heißt Fleisch des Lob-Opfers, das ist, neben dem Lob-Opfer; item, Lob-Opfer des Dank-Opfers.
\quellseite{188}

\chapter{Das 8. Kapitel}

\section{Einweihung der Priester}
% Omitted: numbered Bible verse, excluded; printed folio or catchword marker; running head; not text apparatus.
\quellseite{190}

\chapter{Das 9. Capitel}

\section{Das erste Opfer Aarons wird vom Feuer verzehret}
\quellseite{191}

\chapter{Das 10. Kapitel}

\section{Nadab und Abihu vom Feuer getödtet}

\section{Summarium}

Ist ein schrecklich Exempel, und zur Warnung fürgeschrieben, auf daß jedermann sich hüte und nicht nach seinem Gutdünken, noch eigener Andacht, sondern allein nach dem Wort und Befehl Gottes handle, in den Sachen, die Gottes Ehre und den rechten Gottesdienst betreffen. Denn da die Söhne Aarons solches nicht thäten, mußten sie darum sterben, V. 2.

a Ich werde geheiliget) Qui mihi ministrant, sciant me esse sanctum, non prophanum. Das ist: die mir dienen, sollen wissen, ich sei heilig, nicht unheilig.

b Sein Blut) weil sein Blut nicht ins Heilige gebracht ist, sollte es nicht verbrannt, sondern gegessen worden sein. Welches Blut aber hinein gebracht wär, mußte man nicht essen, sondern alles verbrennen, v. 4. 5. c. 6, 30. und in c. 16. 26. 27.
\quellseite{192}
Darnach ists eine feine Warnung, da Mose Aaron und seinen Söhnen gebeut, daß sie nicht sollen Wein noch stark Getränke trinken, wenn sie in die Hütten des Zeugnißes gehen, auf daß sie nicht sterben, V. 8. Also vermahnet auch Christus seine Christen, daß sie sollen nüchtern sein. Luc. 19, 34. Und sonderlich St. Paulus 1 Tim. 3, 3. lehret: Ein Bischof soll nicht ein Weinsäufer sein. Denn solche Unart soll an keinem Christen und sonderlich an keinem Kirchendiener sein.

\section{Votum}

Der gnädige und treue GOtt regiere uns alle, daß wir in Gehorsam seiner Gebote wandeln, und daß wir uns für allen Sünden hüten, die zu Unehren seinem heiligen Namen, und uns zu Schand und Schaden gelangen mögen, Amen.

\chapter{Das 11. Kapitel}

\section{Vorrede}

Wie wir Gott für seine vielfältigen großen Wohlthaten selten herzlich danken: also gebrauchen wir solche selten recht. Gott hat uns durch JESUM Christum seinen lieben Sohn und unsern treuen Heiland befreiet von der Beschwerung der levitischen Gesetze Mosi und andern; wer danket Gott dafür? Also mißbraucht fast jedermann diese Freiheit, daß der Sünden alle Winkel der Welt voll sind. Derowegen laßt uns Gott danken, daß wir Freie sind, und laßt uns leben, nicht als hätten wir die Freiheit zum Deckel der Bosheit, sondern als die Knechte Gottes, 1 Pet. 2. Denn wir sind nicht zur christlichen Freiheit gekommen, daß uns frei stehe zu thun und zu lassen, was wir wollen, und also die Freiheit mit dem Evangelio zum Schand-Deckel der Sünden zu gebrauchen; sondern daß wir Gott dienen sollen, wie er uns solches in seinem Gesetz und Wort befohlen hat. Davon sagt St. Paulus Röm. 6. also: Nun ihr frei worden seid von der Sünde, seid ihr Knechte worden der Gerechtigkeit. Dies unterrichtet auch das 11. Kapitel.

\section{Vom Unterschied reiner und unreiner Thiere}

a Diese vier Thiere sind in unsern Landen nicht, wiewohl gemeiniglich Arbe und Hagab für Heuschrecken gehalten werden, die auch vierfüßige Vögel sind. Aber es ist gewisser, diese hebräische Namen zu gebrauchen, wie wir mit Halleluja und andern fremden Sprachen Namen thun.
\quellseite{193}

\section{Summarium}

Dies Kapitel sammt etlichen andern, so hernach folgen, sind von etlichen sonderlichen Gebräuchen, welche die Juden aus Gottes Ordnung und Befehl halten mußten, als mit dem Essen, daß ihnen etliche Speise verboten war; item mit der Unreinigkeit am Leibe, daß sie nicht durften unter die Leute gehen etc. Solches sind eitel Beschwerung gewesen, damit Gott dies Volk nicht allein hat wollen unterscheiden von den Heiden, sondern auch äußerlich beschweren, daß sie zum ersten genugsam Ursache hätten, den Gehorsam gegen Gott zu erzeigen; darnach daß sie nicht Raum noch Weile hätten, ihnen eigene Gottesdienste zu erwählen.

Wir Christen aber sollten Gott von Herzen danken, daß wir durch Christum von solchen äußerlichen Beschwerungen erlöset sind; und nicht allein solches zu halten nicht mehr schuldig, sondern auch von der höchsten Last des Gesetzes gefreiet sind, nämlich vom Fluch und Tode, so um der Sünden willen auf uns liegt.

\section{Votum}

GOtt sey Lob, Ehr und Preis für alle seine geistliche und leibliche Wohlthaten;
\quellseite{194}
der verleihe, daß mir durch diese seine Wohlthaten im Glauben gestärket werden, und unser Leben derowegen wohl bessern, Amen.

\chapter{Das 12. Kapitel}

\section{Vorrede}

Der Kinder Geburt ist ein unaussprechlich groß Mirakel und Wunderwerk. Aber wie wenig und selten danken Gott viele Menschen, daß er sie im Mutterleibe erschaffen und erhalten hat, und daß er sie auch mit vollkommenen Gliedmaßen, frisch und gesund in diese Welt hat lassen geboren werden. Derowegen wollen wir solche und andere Gottes Wohlthaten erkennen, und ihm alle Tage dafür danken, und darum ferner anrufen. Denn die Gesundheit des Leibes und der Glieder desselben nicht sein ein natürlich Gut, oder das ungefähr entstehet, wie viel solche Menschen es dafür halten; sondern sind eine besondere edle Gabe Gottes. Und schreiben die Anatomici und gelehrten Aerzte, daß der menschliche Leib so viel Glieder habe, als Tage im Jahre sind. Sollen wir denn nun nicht Gott danken, und bitten täglich für die Gesundheit unsers Leibes und unserer Glieder? Dies erinnert auch das 12. Kapitel.

\section{Ordnung der Kindbetterinnen}

\section{Summarium}

Das ist eine mosaische Ordnung von den Sechswöchnerinnen oder Kindbetterinnen, welcher wir im Neuen Testament eben so wenig verbunden sind, als den andern von Essen und unreiner Speise. - Darum ists ein Mißbrauch und unwesen, daß man das Einsegnen im Papstthum, nach den Sechswochen oder Kindbett hat für nöthig erachtet.

Daß man aber die Sechswöchner halten, geschieht nicht um Moses Ordnung willen, sondern daß es die Natur und Noth also erfordert, daß eine Kindbetterin nach der Geburt sich innen halten und eine Zeit lang ruhig sein, und mit gebührlichem Essen und Trinken ihr auswarten soll; nicht allein ihret halben, daß sie wieder zu Kräften komme, sondern auch des Kindes halben, daß das nicht verwahrloset werde, weil es noch schwach ist, und seine Nahrung von der Mutter haben muß.

\section{Votum}

Wir danken Gott herzlich für alle Wohlthaten, die er uns durch unsere liebe Eltern erzeiget hat; der segne und behüte alle schwangere Hausmütter und Sechswöchnerinnen, mit ihren Kinderlein, und verleihe, daß wir mit ihnen allen selig werden, Amen.

\chapter{Das 13. Kapitel}

\section{Kennzeichen des Aussatzes an den Menschen und Kleidern}
\quellseite{195}

a Die ganze Haut) dieser Aussatz heißt rein, denn es ist ein gesunder Leib, der sich also selbst reiniget, als mit Pocken, Masern und Krätze geschiehet, da durch den ganzen Leib das Böse heraus schlägt; wie wir Deutschen sagen, es sei gesund etc.
% Omitted: numbered Bible verse, excluded; printed folio or catchword marker; running head; not text apparatus.
\quellseite{197}

\chapter{Das 14. Kapitel}

\section{Reinigung des Aussatzes}

a Log, Log ist ein klein Maßlein, auf Hebräisch also genannt, aber noch ungewiß, wie groß es sei.
% Omitted: Bible narrative, running matter, or non-apparatus block; numbered Bible verse, excluded; printed folio or catchword marker; running head; not text apparatus.
\quellseite{199}

\chapter{Das 15. Kapitel}

\section[Manns- und Weibspersonen, mit unreinem Fluß behaftet, wie]{Manns- und Weibspersonen, mit unreinem Fluß behaftet, wie sie zu reinigen}
\quellseite{200}

\chapter{Das 16. Kapitel}

\section{Jährliches Versöhnopfer}
\quellseite{201}

a Unrein) das ist, ganz Israel ist unrein und in Sünden, und Gott rechnet noch mitten unter viel Sündern, und heiliget sie; das ist, keine Heiligkeit genießen sie, und heißen heilig um seinetwillen, sonst ist alles unrein und Sünde mit ihnen, das ist Christus in seinem Volk.

\chapter{Das 17. Kapitel}

\section{Ort der Opfer bestimmt. Blut zu essen verboten}

a Schuldig) da siehst du, daß er \unsicher{nicht will einen selbst erwählten Gottesdienst} haben außer seinem geordneten Gottesdienst. Darum nennet er einen solchen Opferer einen Mörder, wie er 1 Mos. 9, 6. auch thut.
\quellseite{202}

b Priester nicht sie selbst, aus eigener Wahl und Andacht.

\chapter{Das 18. Kapitel}

\section{Vorrede}

Wie wir uns für andern Sünden hüten sollen: also die, so sich in Ehestand begeben wollen, sollen \unsicher{Ehetrei} fliehen und meiden, das ist, sie sollen nicht zu nahe in die Blutsfreundschaft oder Schwägerschaft freien. Denn Unzucht treiben ist große Sünde, die Ehe brechen ist größer; in Blutschanden leben ist die größte. Aber da sind viel evangelische, eigenwillige Leute, die keiner guten Ordnung, ja auch Gottes Gebot nicht gehorsamen, und ziehen also Blutschulden und Gottes Strafen auf Land und Leute. Wie denn dieser Sünden Mose gedenkt als der fürnehmsten Ursachen einer der Sündflut, 1 Mos. 6. Solches zu verhüten, soll jedermann bekannt sein das 18. Kapitel des dritten Buches Mose.

\section{Verbotene Grade der Blutsfreundschaft im Heirathen}
\quellseite{203}

a Moloch) war ein Abgott, dem sie ihre eigenen Kinder zu Dienst verbrannten; wie Manasse that, der König Juda, und meineten, Gott damit zu dienen, wie Abraham, da er Isaak, seinen Sohn, opfert. Aber weil ihnen Gott das nicht befohlen hatte, wie dem Abraham, war es unrecht. Darum spricht hie Gott, daß sein Name dadurch entheiliget werde; denn es geschah unter Gottes Namen und war doch teuflisch. Wie auch jetzt die Kloster-Gelübde und andere Menschen-Aufsätze viele Leute verderben, unter göttlichem Namen, als sei es Gottesdienst.

\section{Summarium}

Das ist eine Ordnung, die da belanget das sechste Gebot, in welcher angezeigt wird, welche Personen zum ehelichen Beiwohnen erlaubt oder verboten sind.

Solche Ordnung ist nicht zu achten wie andere Gesetze, so von Mose geboten, nun aber im Neuen Testament aufgehoben sind; sondern sie zeucht sich auf die Natur. Denn der Text meldet klar, wie solches unordentliches Beiwohnen auch unter den Heiden eine Sünde sei, V. 24. Und daß um solcher Sünden willen, als Ehebruch und andere Unzucht, Land und Leute ausgeräumet und verwüstet müssen werden, V. 25. Darum soll jedermann sich für Unzucht hüten, und sonderlich der Blut-Freundschaft schonen, weil es Gott so theuer verboten und so greulich auch unter den Heiden gestraft hat.

\section{Votum}

Helfe der treue GOtt, daß wir seinen heiligen Geboten folgen, und als Gehorsame an Leib und Seel von ihm gesegnet werden, um JEsu Christi seines lieben Sohns Gehorsam und Verdiensts willen, Amen.

\chapter{Das 19. Kapitel}

\section{Vorrede}

Wenn zusammen sollte gerechnet werden, wie oft die meisten Menschen in der ganzen Welt an die heiligen zehn Gebote gedenken, oder wie
\quellseite{204}
oft sie dieselbigen bei ihrem Gebet sprechen: ach, die Summa würde klein und geringe sein. Aber wie Gott durch Mosen oft die heiligen zehen Gebote hat wiederholen und erklären lassen: also will Gott, daß alle Menschen die heiligen zehen Gebote täglich sprechen und betrachten sollen, daß sie dadurch zur Buße gereizt werden, daß ist, daß sie ihre Sünde erkennen, die Gott abbitten, des Herrn Christi im wahren Glauben sich trösten, und darauf ein Gott wohlgefälliges und ihrem Nächsten nützliches Leben führen. Darauf höret davon an das 19. Kapitel.

\section{Auslegung der zehn Gebote, sammt andern Gesetzen}

a) Gefallen könnte) das ist: Macht's, wie es gebeut, nicht wie euch dünket oder gefällt nach eigener Andacht. Denn er will schlechthin keine selbst erwählte noch eigen erdachte Weise haben. Darum sehet zu, daß ihr also opfert, daß ihm gefalle, und nicht weiter erzürnet werde durch eure eigene Weise. Nam praeter juge sacrificium non erat necessarium vovere et offerre ei praeterea pro peccatis; illa pacifica sint libera. Das ist: denn außer dem täglichen Opfer war nicht nöthig, Gelübde und Opfer zu thun überdieß noch für die Sünden; jene Dankopfer sollen freiwillig sein.
\quellseite{205}

b Vorhautschneiden ist hie so viel, als drei Jahr barren, wie er selbst deutet und spricht: Drei Jahr sollt ihr sie achten für unbeschnitten etc.

\section{Summarium}

Dies ist eine Predigt, in welcher Gott etliche Stücke aus den zehen Geboten wiederholet und ausstreicht, nämlich, wie man den Nächsten lieben, ihm dienen und allerlei Freundschaft beweisen soll. Item, wie man sich für Aberglauben, Wahrsagen und dergleichen Sünden hüten soll, V. 31.

An solche Gebote hänget er alle Wege diesen Beschluß: Ich bin der HErr euer Gott; damit er anzeiget, daß wir Gott Gehorsam schuldig sind, in allem, was er haben will, es sei so gering es wolle. Auf daß wir allein auf sein Wort und Befehl sehen, und uns nichts an solchem Gehorsam hindern lassen.

\section{Votum}

GOtt stärke unsere schwachen und furchtsamen Herzen, und verleihe, daß wir in Noth und Tod seinen Geboten gehorsam sein, und in wahrem Glauben, Liebe, Hoffnung und Geduld leben und sterben. Amen.

\chapter{Das 20. Kapitel}

\section{Vorrede}

Gott, als unser Schöpfer und lieber Vater, meinet ja treulich mit uns Menschen, auf daß er uns segnen könne und nicht strafen dürfe. Denn beides thut er in seinem Wort: er lehret, für was Sünden wir uns in allen Ständen hüten sollen, und zeiget an, daß er die Sünden nicht wolle ungestraft lassen, sondern die an allen Unbußfertigen mit zeitlicher und ewiger Plage und Marter belohnen. Derowegen, wer Gottes Treue erkennet, der lebe täglich nach seinem (im Wort) offenbarten Willen; und wer ihm selbst Gutes gönnet, der hüte sich für allen Sünden, so hütet er sich auch also vor den Strafen der Sünden. Sonst heißt es, wie Gott saget in diesem dritten Buch Mose am 26. Kapitel: Werdet ihr mir nicht gehorchen und nicht thun diese Gebote alle, und werdet meine Satzungen verachten, und eure Seele meine Rechte verwerfen, daß ihr nicht thut alle meine Gebote, so will ich euch auch solches thun: ich will euch heimsuchen mit Schrecken, Schwulst und Fieber, daß euch die Angesichter verfallen und der Leib verschmachten soll. Darauf höret ferner an das 20. Kapitel.

\section{Strafen unterschiedlicher Sünden}

a Das ist, seiner Kinder.
% Omitted: numbered Bible verse, excluded; printed folio or catchword marker; running head; not text apparatus.
\quellseite{207}

\section{Summarium}

Die ist eine Predigt, in welcher Gott etliche Gesetze wiederholet, wie man den Abgöttern nicht opfern, Wahrsagen und Zeichendeuten nicht brauchen soll, Gott wolle sonst selbst, wo es schon die Leute nicht thäten, solche Sünde strafen, V. 4. 5.

Sonderlich aber ist zu merken, daß er befiehlt, mit dem Tod zu strafen, wo jemand den Eltern fluchet, V. 9., seine Ehe bricht, V. 10., oder unnatürliche Unzucht treibet, V. 15. Item, daß er die, so zu nahe in die Blutsfreundschaft sich verheirathen, mit Unfruchtbarkeit strafen und ohne Kinder wolle sterben lassen, V. 20. 21. Solche Gebote ziehen sich fast alle auf die Natur, darum binden sie uns gleich so wohl als die Juden.

\section{Votum}

Der treue barmherzige GOtt, der da befiehlet, was wir thun sollen, der verleihe Gnade und Hülfe durch den heiligen Geist, daß wir auch seinem Befehl folgen und nachkommen mögen, und in seinem Gehorsam leben und sterben, Amen.

\chapter{Das 21. Kapitel}

\section{Wie sich ein Priester zu verhalten}

a Krone) auf Hebräisch Nezer, hiervon die Nazarei heißen, daß sie sich enthielten, und sonderten von weltlicher Speise, 4 Mos. 6, 3. Und Simson von Mutterleib ein Nazareus hießet, Richt. 13, 5. Und unser HErr Christus auch daher der rechte Nazareus heißet von Nazareth, Matth. 2, 23.; als der von aller Sünde rein und heilig und kein Unheiliges noch Ungeweihtes an ihm ist. Wiewohl die Juden auf Neid und Feinheit ihn nicht Nazar, sondern Nozri, das ist, den Verächter und Verstörten, wie einen Schänder nennen, so sie doch selbst die rechten Nazirim sind, in alle Welt zerstreuet und verstöret. Und mich dünkt, daß sich St. Paulus Röm. 1, 1. einen Nazareum nennet, da er sich einen Abgesonderten rühmet zum Evangelio, vom Gesetz c. 1 Tim. 3, 2. Tit. 1, 7.
\quellseite{208}

b (ungewöhnlich und seltsam) das sind allerlei Ungestalten. Als so die Nasen zu groß, zu klein, zu krumm, zu breit, item, das Maul krumm, schartig, blöcket \&c., das einen für Andern verstellet.

\chapter{Das 22. Kapitel}

\section{Von des Opfers Beschaffenheit}

a (Enthalten) das ist, wie die Nazaräer sich heiliglich abhalten und meinen, Nazaräer etc., nämlich, daß sie nicht essen sollen vom Opfer, sie seien denn ganz rein und heilig.
\quellseite{209}

b Fehl) über dies Stück klagt Maleachi 1, 8. sehr hart; denn was nicht taugt noch gut ist, das gibt man Gott und seinen Dienern, und wills also ertragen. Er aber doch solches hie als unangenehm verdammt, denn er hat lieb einen fröhlichen Geber.

c Wandelbar) als da nur ein Ohr oder keins, das ein Auge oder keins, daß ein oder mehrere Füße mangelt, oder sonst ungestaltet und ungeschaffen ist; es heißt: Wer geben will, der gebe was Gutes, oder lasse es anstehen.

\chapter{Das 23. Kapitel}

\section{Vorrede}

Den Weltkindern sind alle Tage Sünden-Tage; dagegen sind den Christen alle Tage Sabbath-Tage, das ist, solche Tage, an denen sie, wo sie nur können, Gottes Wort hören, an denen sie auch Gottes Wohltaten betrachten, ihm dafür danken, für sich, die Ihren, und für die Noth der ganzen Christenheit beten, von Sünden ruhen und die Tage mit heiligen Gott wohlgefälligen Werken, wie es ihr Beruf und Stand erfordert, zubringen. Derowegen, weil wir auch Christen sein, so wollen wir dergleichen alle Tage, als Sabbaths-Tage, heilig halten und christlich begehen. Dazu uns denn helfe unser lieber HErr JEsus Christus, hochgelobet in Ewigkeit, Amen. Darauf höret ferner an das 23. Capitel.

\section{Ordnung der vornehmsten Feste}

a Dienstbarkeit) als da sind die Werke, so man an den Werktagen thut, Nahrung zu suchen, da man Getreide und Vieh zu brauchet. Aber Haus-Geschäfte und Gottesdienst ist nicht verboten, als kochen, kehren, kleiden etc.
\quellseite{210}

b Blasens zum Gedächtnis] solch Blasen mit einem Horn geschah, daß man damit Gottes und seiner Wunder gedachte; wie er sie erlöset hatte, davon predigte und dankte. Wie bei uns durchs Evangelium Christi und seiner Erlösung gedacht und geprediget wird.
\quellseite{211}

\section{Summarium}

Hie ist eine Ordnung von den fürnehmsten Festen, so Gott durchs ganze Jahr den Juden hatte befohlen zu halten. Welche Feste alle mit einander dahin gestiftet sind, auf daß man Gottes Güte und Wunderwerke nicht aus dem Gedächtnis lasse, sondern Gott dafür soll dankbar sein, und sich seiner Hülfe auch im künftigen Unglück trösten.

Den Sabbath nennt er das fürnehmste Fest, V. 3. Denn allewege der siebente Tag dahin gestiftet ist, daß man Gottes Wort hören und Gottes Werk betrachten und Gott für einen Schöpfer Himmels und der Erden erkennen und loben soll, von welchem alle Menschen Leib und Leben haben und alles gewarten, bitten und hoffen sollen, was zur Erhaltung Leibes und Lebens und endlich zur Seligkeit vonnöthen ist. Solches heißt den Sabbath nicht allein feiern, sondern auch heiligen. Es ist aber solches Fest oder Sabbathtag nicht allererst von Mose geordnet, sondern bald nach Vollendung der Kreaturen, ehe denn Adam und Eva sind gefallen, von Gott selbst aufgesetzt und zu solchem Gottesdienst geheiliget, das ist, dazu verordnet und darauf geordnet, daß man ihn mit heiligen Dingen soll zubringen. Solche Satzung hat Moses verneuert. Aber da sind der Propheten Zeugniß und Predigt immer daneben gegangen, im Neuen Testament: wenn der verheißene Same kommen werde, werde solcher mosaische Sabbath fallen und alle Tage Sabbath sein, für und für, wie Christus auch sagt: Solches thut zu meinem Gedächtnis, so oft ihrs thut. Und Paulus: Verkündiget den Tod des HErrn, bis er kömmt, das ist, alle Tage für und für rühmet Gottes Gnade und Wohlthat in Christo. 1 Cor. 11, 26.

1. Darnach waren andere sonderliche Feste. Das erste war das Osterfest, V. 6. Daran man bedenken sollte, wie Gott sein Volk wunderbarlich aus Egypten geführet und den Tyrannen Pharao mit seinem Volk gestürzet hat, und Gott für solche Erlösung dankbar sein, und sich in anderem Unglück seiner Hülfe auch trösten. Neben solchem Dank war dies Fest auch eine Erinnerung der künftigen Erlösung, durch das rechte Osterlamm, den Sohn Gottes, unsern HErrn Jesum Christum, der nicht allein die Juden, sondern alle Welt von dem schädlichen Pharao und Egypten, das ist, von Sünde, Tod und Hölle erlösen sollte. Wie St. Paulus solches selbst deutet 1 Cor. 5. V. 7.

2. Das andere war das Fest der Erstlinge, V. 10., daß allewege, wenn der Schnitt anging, man eine Garbe von einem jeden Gewächs im Tempel opfern mußte, sammt einem Brandopfer, Speisopfer und Trankopfer; und darnach fünfzig Tage zählen, bis die Feld-Arbeit gar verrichtet ward. Da mußte aus einem jeden Hause Brod von neuem Getreide in dem Tempel geopfert werden. Solches geschah, auf daß die Leute lernten Gott für seinen Segen danken und erkennen, daß es Gottes Gabe sei, wenn der Feld-Bau wohl geräth, und Gott allerlei, was wir bedürfen, reichlich läßt wachsen. Man soll aber solchen Segen nicht allein dahin brauchen, daß wir alles vollauf und genug haben, sondern fürnehmlich damit helfen den äußerlichen Gottesdienst erhalten und den Kirchendienern ihren Theil davon geben, auf daß sie sich desto füglich erhalten, ihres Amts warten und zur Erkenntniß Gottes und Seligkeit Vielen helfen können. Solcher Ursache halben hat Gott dies Fest geordnet.

3. Darnach war das Posaunenfest, V. 24., welches man hielt am ersten Tage des siebenten Monats nach Ostern, welches bei uns ungefährlich ist der Weinmonat, wenn alle Frucht des Feldes, Korn, Wein, Oel und was da wächst,
\quellseite{212}
eingeführet sind. Solches Fest scheinet auch, daß es am fürnehmsten dahin hat sollen dienen, auf daß man Gott für seinen Segen dankbar wäre.

4. Am zehenten Tage nach solchem Fest war das Fest der Versühnung, B. 27., da mußten die Juden fasten und feiren. Denn es gehört sich, wenn Gott seinen Segen gibt, daß wir erkennen, wie wirs nicht werth sind und wohl alles Unglück unserer Sünden halben verdienet hätten, wo Gott sonst nicht so gnädig wäre. Darum mußte man am selbigen Feste die Sünde bekennen und Gott um Versühnung und Gnade bitten. Aus diesem Fest haben die Papisten die Weise genommen, daß sie zur Genugthuung für die Sünde den Leuten sonderliche Festtage haben aufgelegt. Aber es ist solches Fasten schlechter Gottesdienst, ja auch wohl ein Greuel für Gott; sintemal Gott davon uns nichts geboten, sondern Vergebung der Sünden im Namen Jesu Christi allein zu predigen befohlen hat, und uns gar nichts angehet, was Gott in solchem Fall den Juden befohlen hat. Was aber das für eine Versühnung gewesen, ist im Anfang dieses Buches gemeldet.

5. Das letzte Fest war am fünften Tage nach dem Versühnfest und hieß das Fest der Laubhütten, B. 34. Da mußten die Juden unter keinem Dach, noch in ihren Häusern weder essen, schlafen, noch anders thun, sondern mußten unter dem offenen Himmel bleiben sieben ganze Tage, zum Gedächtniß, daß sie Gott in der Wüste ohne solchen Segen, Weins, Korns und anderes Gewächs wunderbarlich erhalten und ernähret hat.

Auf daß wir auch lernen, obgleich solche Feste mit andern Moses-Ordnungen gefallen sind, daß wir dennoch Gott um alle seine Wohlthaten dankbar sein und seinen Segen gern zu Unterhaltung der Kirchendiener und des rechten Gottesdienstes brauchen sollen. Sonderlich aber haben wir im Neuen Testament, wie oben auch gemeldet, alle Tage unsern Sabbath, unsere Ostern- und Versühn-Feste, weil Christus für uns gestorben ist und sein Verdienst durch sein Wort uns geschenket und zu eigen gibt. Für solche Wohlthat, neben dem zeitlichen Segen, sollen wir Gott alle Tage (weil wir hie gleich wie die Juden in der Wüste, ehe sie in das gelobte Land und die rechte Heimath kommen, in unserer Pilgerschaft sind) danken und ihn loben, bis wir endlich gar selig werden und in die ewige Hütte kommen, Amen.

\chapter{Das 24. Kapitel}

\section{Vorrede}

Wenn sonst keine Sünde auf Erden wäre, denn die große Gotteslästerung und der gräuliche Mißbrauch göttlichen Namens, so wäre es Sünde genug, Gottes Zorn und alle Strafen zu erwecken und auf die Welt zu ziehen; und wäre freilich nicht Wunder, daß Gott nicht allein mit Pestilen, Krieg und Theurung die Welt strafe, sondern auch wohl mit Schwefel und höllischem Feuer auf sie schienen und schlacken ließ. Derowegen wer an dieser Sünde schuldig ist, der thue heute Buße, und wer dieser Sünde wehren kann, der thue es bald, um Gottes willen. Denn versucht sei die Zunge, die lästerlich redet wider den, der sie geschaffen hat, sagt Gordius der Märtyrer. Dagegen mag man auch sagen: Selig sei die Hand, welche die Gotteslästerung strafet, und die es wehren soll und kann. Davon höret auch an das 24. Kapitel.

\section{Vom Leuchter. Schaubrode. Strafe der Gotteslästerer und Todtschläger}

a Denkbrode) das sind die Schaubrode, welche die Kuchen heißen, darum, daß sie breit waren wie Kuchen. Und sind Denkbrode, darum, daß sie damit Gottes gedenken, und von ihm predigen sollen. Gleich wie Christus uns befehlet, daß wir sein gedenken, das ist, seinen Tod verkündigen und predigen sollen, 1 Cor. 11.

b Feuer) das ist Opfer.
\quellseite{213}

\section{Summarium}

Am ersten ist da ein Exempel, wie das Volk von dem Thron das Licht im Tempel erhalten und Oel dazu habe geben müssen, B. 2. Denn Gott will, daß jedermann etwas dazu geben soll, daß rechter Gottesdienst erhalten werde.

Sonderlich aber reimet sich die geistliche Deutung fein. Das Licht im Tempel ist das Predigtamt und die Kirchendiener, wie Christus seine Apostel Matth. 5, 14. auch nennt, das Licht der Welt. Solche sollen ihren Unterhalt von der Gemeine haben, wie die Juden zum Licht das Oel mußten geben.

Darnach ist ein ernst Exempel, wie Gott gebeut den zu steinigen, der bei Gottes Namen geflucht hatte, V. 13., auf daß wir Gottes Namen, Gott zu Lob und Ehre, lernen recht brauchen, mit Predigen, Beten und Bekennen, und nicht also zu Unehren mißbrauchen, wie leider jetzt in der Welt geschiehet und wir derhalben erstlich alle insgemein, und darnach ein jeder in seinem Hause mancherlei Unglück müssen tragen. Denn es ja eine gräuliche, schreckliche Sünde ist, daß wir aus schändlicher böser Gewohnheit unsers lieben Herrn Christi, seiner Taufe, Sacrament, Leiden und Marter, welches alles um unsrer Seligkeit willen geschehen, dahin mißbrauchen, daß wir andern Böses dadurch wünschen. Und ist gewiß, solche Sünde verdammt Leib und Seele an vielen, daß sie hie all ihr Glück und Segen durch solches Fluchen und Lästern selbst verschlagen und ihnen das höllische Feuer verdienen. Darum lerne ein jeder Christ dafür sich hüten, sonderlich aber soll weltliche Obrigkeit mit harter Strafe solchem Aergerniß wehren und Gott seine Ehre retten, so viel möglich.

\section{Votum}

GOtt regiere die Obrigkeit, daß sie die Gotteslästerung ernstlich strafe, und regiere uns, daß wir seinen heiligen Namen nicht schänden, sondern denselben im Herzen und mit Worten und Thaten loben und preisen, Amen.

\chapter{Das 25. Kapitel}

\section{Vorrede}

Daß jetziger Zeit allerlei Händel und Gewerbe stecken bleiben, dessen ist nicht eine kleine Ursache, daß man allerlei Betrug in Käufen und Verkaufen gebraucht; und, da Gott segnet, daß der meiste Theil so wenig der Armuth davon mittheilet. Derowegen, ob schon die Zeit schwer ist, dennoch nähre sich ein jeder mit gutem Gewissen, und lasse ihm arme Leute befohlen sein. Also wird ihm an Gottes Segen zu täglichem Brod und an Trost, unter gemeinem Kummer, nicht mangeln. Von diesem sagt Salomon in seinen Sprüchen im 19. Kapitel also: Wer sich des Armen erbarmet, der leihet dem HErrn, der wird ihm wieder Gutes vergelten. Von jenem aber spricht also David im 128. Psalm: Du wirst dich nähren deiner Hände Arbeit; wohl dir, du hast es gut. Dies unterrichtet auch das 25. Kapitel.

\section{Feier- und Jubeljahr}
\quellseite{214}

a Speise) das ist, gemein sein und nicht einsammeln noch ausschütten.
\quellseite{215}

b Ueberfaß) Wucher heißt er, so mit Geld geschiehet; Ueberfaß, wenn der arme Mann muß kaufen oder annehmen die tägliche Waare, so theuer der Geizhalß will, weil ers haben muß zur Noth.

\section{Summarium}

Hier ist eine Ordnung etlicher äußerlicher Händel, welche alle dahin gerichtet sind, daß der Armen nicht vergessen und sie mit Kaufen und Verkaufen nicht übervortheilet, noch beschweret werden. Denn daß bei den Juden allewegen das siebente Jahr das Feld ganz und gar feiern mußte, geschah, wie der Text meldet, v. 6., auf daß alles, was von sich selbst wüchse, jedermann gemein wäre und die armen Leute sammt dem Vieh ihre Nahrung desto besser haben möchten.

Digitized by Google
\quellseite{216}
Also ist die Ordnung mit dem fünfzigsten Jahr auch gewesen: obschon die Leute verdarben und ihre Aecker und Güter Armuth halber verkaufen mußten, daß sie doch im Jubel-Jahr wieder zu Gütern kamen, V. 10. Solches ist eine überaus gute und gleichmäßige Ordnung gewesen, daß die Güter nicht haben können ganz verruckt werden.

Sonderlich aber ists fein, daß Gott will, daß in Kaufen und Verkaufen allein die Billigkeit und kein beschwerlicher Vortheil soll gesucht werden, V. 16. Und gehet, man soll die Armen mit Wucher und Aufsatz nicht beschweren und sonst auch allewege einen Bruder gnädiger halten, denn einen Fremden, V. 25. 35. 39. 47.

\section{Votum}

Helfe der treue GOtt, daß diese Lehren in uns Nutz schaffen, und daß ein jedes also lebe, auf daß es an Sünden arm, und an Tugend reich sei und bleibe, und daß wir alle um Christi willen das ewige Jubeljahr im Himmel halten mögen, Amen.

\chapter{Das 26. Capitel}

\section{Vorrede}

Gott will, daß kein Mensch verloren werde, sondern daß sich ein jeder bekehre, seiner Barmherzigkeit und Güte genieße und selig werde. Derowegen locket er zur Buße durch ernste Predigten und Strafen. Solche Treue und Gnade Gottes lasset uns erkennen und wohl gebrauchen; daß ein jedes von seinen Sünden ablasse und dagegen thue, was es Gott zu Ehren, seinem Nächsten zu Nutz und ihm selbst zur Wohlfahrt, nach Gottes Wort, thun soll und kann. Denn je treulicher Gott warnet, nimmt man die Warnung nicht an und führet fort in Sünde, Leben und Unbußfertigkeit, je härter und grausamer er darauf strafet, wie solches die heilige Schrift und die Erfahrung bezeuget. Davon höret auch an das 26. Capitel.

\section{Gedrohter Fluch, und verheißener Segen}
\quellseite{217}

\section{Summarium}

Ist eine herrliche und ernste Predigt, in welcher Gott sein Volk vermahnet, V. 3. Weil sie nun wissen, was sie thun und lassen, wie sie sich mit Opfern und dem Gottesdienst halten sollen, daß sie bei solchem Befehl bleiben und nichts dran ändern, noch unterlassen sollen. Und verheißet ihnen, wo sie solches thun, er wolle mit reichem Segen auf dem Felde, mit Frieden im Lande und mit allem Glück sie überschütten. Wiederum, wo sie es nicht thun, da soll Krankheit, Hunger, Krieg, schädliche Thiere und alles Unglück sie treffen ohne Aufhören, V. 14.

Da sehet ihr nun, wo solcher Jammer in der Welt alle herkömmt; darum mag sich jedermann wohl bessern und frömmer werden, so wir anders solcher Strafen wollen überhoben sein.

a Ueber den andern) wie es geschieht in der Flucht, Furcht und Schrecken.

b Gefallen) das ist: gleich wie sie Lust an ihren Sünden und Ekel an meinen Rechten hatten: also werden sie wiederum Lust und Gefallen haben an der Strafe, und sagen: ach wie recht ist uns geschehen; Dank hab unsere verfluchte Sünde, daß haben wir nun davon. O recht lieber Gott, o recht. Und das sind Gedanken und Worte einer Erkenntniß und Buße, die sich selbst aus Herzensgrund hassen und anklagen lehrt, \unsicher{wozu sich, was hab ich gethan?} daß gefällt denn Gott, daß er wieder gnädig wird. Darum haben wir das Wort Missethat verbüßt die Strafe der Missethat, solchen Verstand zu geben, sonst lautet's, als sollten sie Gefallen an der Missethat haben. Eben so ist's auch zu verstehen: dem Lande gefällt seine Feier, das ist, es spricht: Gott habe Recht in der Strafe, daß es wüste liegen muß um des Volks willen, nach dem es sich sehnet. Solcher Weise redet auch Jes. 40, 2. Dimissa est iniquitas, id est, placita et accepta poena.
\quellseite{218}

\section{Summarium}

Es ist eine heftige und schreckliche Predigt, darin Gott dräuet den unbußfertigen Sündern, daß Gott mit der Strafe nicht aufhören, sondern je mehr man mit Sünden fortfähret, je mehr er mit der Strafe anhalten wolle. Auf daß man lerne, daß Gottes Zorn mit nichts mag abgelegt werden, es sei denn, daß man sich bessere und von Sünden ablasse, und Gnade begehre.

Am Ende ist ein sehr tröstlich Stück, daß Gott solche harte Ruthen über die Sünder nicht darum gehen lasse, daß er sie gar verderben, sondern allein, daß er sie zur Buße treiben und frömmer machen und wieder zu Gnaden annehmen wolle, V. 42. Denn weil Gott anders nicht, denn gnädig ist, will er seinen Gnaden-Bund nicht gar um unserer Sünden willen aufheben. Darum sollen wir in Sünden nicht verzagen, sondern davon ablassen, und uns mit Gottes Barmherzigkeit trösten.

\section{Gebet}

Erleuchte und regiere uns der gütige treue GOtt, daß wir für allen Sünden uns hüten, und dagegen gottselig und ehrbar allezeit leben, auf daß wir der Sünden wegen nicht zeitlich und ewig an Leib und Seele gestraft werden, Amen.

\chapter{Das 27. Kapitel}

\section{Vorrede}

Wir sind alle für Gott Missetäter. Ach eine große schwere Sünde! Denn wir alle den Eid oft und gräulich gebrochen, den wir in der Taufe Gott gethan und da wir zugesagt haben, daß wir wider unser Fleisch und Blut, wider den Teufel und wider die Welt streiten, Sünde meiden und dagegen in vollem Gehorsam der heiligen zehn Gebote Gott dienen wollen. Solches läßt uns erkennen, Buße thun und Gott leisten, durch unser ganzes Leben, was wir ihm schuldig sind. Denn halten wir, was wir Gott verheißen haben, so wird Gott wieder halten, was er zugesagt hat. Darauf höret auch an das 27. und letzte Capitel.

\section{Von Gelübden und Zehnten}
\quellseite{219}

\section{Summarium}

Dies ist eine Ordnung, wie es mit den Gelübden und den Zehnten soll gehalten werden. Weil nun unsere Päpisten viel von ihren Gelübden rühmen, welche ihre Pfaffen, Mönche und Nonnen thun, soll man diesen gewissen und wahren Unterricht davon wissen.

Die Gelübde im Alten Testament sind fast alle dahin gerichtet, daß den Leuten, so zum Kirchendienst verordnet waren, damit Gott gedienet würde, wie man hie aus Mose klar siehet, daß sie die Häuser, Acker, Vieh, ihre eigenen Leiber eine Zeit lang, verlobet haben und bisweilen mit Geld wieder abgelöset, welches Geld alles dem Tempel und Gottesdienst ist zu Nutz kommen, B. 7.

Solche Gelübde haben dreierlei sonderliche Anhänge gehabt. Erstlich ist es den Juden nicht frei erlaubt gewesen zu geloben, was und wohin sie wollten, sondern es ist hie in Mose das Ziel gesetzt, und mit Gottes Befehl ausgedrückt, was und wohin man geloben soll.

Zum andern sind solche Gelübde nicht ewig gewesen, sondern haben nur eine Zeit lang gewähret, und nach der bestimmten Zeit wieder aufgehöret, ohne daß Gott etlichen Personen insonderheit hat auferlegt, ihr Leben lang sich zu halten, daß anderen nur eine Zeit lang geboten, wie man siehet mit Simson und andern.

Zum dritten waren solche Gelübde dermaßen, daß es allwegen in der Oberherrn Hand stund, dieselbigen zu brechen. Kinder und Gesinde konnten nichts geloben, wider Wissen und Be%
willigen\quellseite{220} des Vaters oder Haußherrn. Ein Weib durfte nichts geloben, ohne Begunst und Verwilligung des Mannes, wie man klar in Mose siehet.

Wenn wir nun unserer Papisten Gelübde gegen diese halten, so ist es ganz und gar mit ihnen umgekehret. Denn erstlich sollen nach des Papstes Recht Vater und Mutter nicht Macht haben, daß sie ihrer Kinder und Ehehalter Gelübde brechen mögen, sondern in solchem Fall billiget es der Papst, daß der Kinder Wille fortgehe und die Eltern ihnen müssen weichen.

Zum andern sind sie nicht auf eine gewisse Zeit bestimmt, sondern bleiben ewig.

Zum dritten, und welches das größte ist, geschehen solche Papst-Gelübde ohne Befehl Gottes, aus eigenem Fürnehmen, ja auch wohl wider Gottes Befehl und Ordnung, wie wir hernach hören werden.

Nun, solche Gelübde des Alten Testamentes, welche nach Gottes Wort und Willen geschehen, sind dermaßen gewesen, daß dennoch die Propheten, wie man im Psalter und sonst siehet, immerdar dies Gelübde weit darüber gerühmet haben, daß man Gott preisen, ihn für gnädig und barmherzig halten und um seiner gnädigen Hülfe willen von Herzen danken soll. Wie es denn der heil. Geist ein Gelübde heißt, Ps. 22. V. 26.: Dich will ich preisen in der großen Gemeinde; ich will mein Gelübde bezahlen, für denen, die dich fürchten. Und im 61. Psalm V. 9.: Ich will deinem Namen lobsingen ewiglich, daß ich mein Gelübde bezahle täglich.

Dies Gelübde des Dankens und Lobens setzen die Propheten weit über die andern Gelübde des Gesetzes. So dagegen unsere Papisten heutiges Tages, auf daß sie ihre Gelübde erhalten mögen, dies recht hohe Gelübde mit Füßen treten, Gottes Wort verfolgen, seinen Namen lästern und schänden, wie man für Augen siehet.

Neben dem, daß ihre Gelübde an ihnen selbst unnütz und wider Gottes Ordnung sind. Denn daß sie die Ehe verschwören, geschieht ja wider Gottes Ordnung, welcher Mann und Weib geschaffen und seinen Segen ihnen verheißen hat. Solch Gelübde aber ist darum desto ärger, weil es so viel Aergerniß, Schande und Laster mit bringet, daß es ein groß Wunder ist, daß Gott nicht längst mit allen Klöstern und Stiften ist umgegangen, wie mit Sodom und Gomorra. Wie können sie sich aber rühmen, daß sie Gehorsam schwören, weil sie aus aller ordentlichen und von Gott bestätigten Gewalt sich ziehen, und niemand, denn ihrem Abgott zu Rom, wollen gehorsam sein.

Fürs dritte ist sehr schimpflich, daß sie rühmen das Gelübde der Armuth, weil jtzund die höchste Klage ist, man ziehe ihre Güter ein, nehme ihnen ihr Einkommen. Wenn sie denn Armuth geloben und dieselbige halten, wie reimet es sich denn, daß sie klagen: man nehme ihnen ihre Güter.

Also ist es am Tage, daß unserer Papisten Gelübde wider Gottes Ordnung und Befehl, ärgerliche und erdichtete Gelübde sind. Darum soll niemand dieselbigen sich binden lassen, sondern eine freie offene Hand haben, dieselbigen zu brechen, ja wissen, wo mans nicht brechen, sondern halten will, daß man wider Gott, sein Wort und Ehre thut. Sintemal solche Gelübde ohne allen Gottes-Befehl geschehen, und wir wissen, daß man Gott mit allen Menschen-Satzungen nicht kann noch soll dienen.

Im Alten Testament waren die Gelübde recht. Ursach, sie geschahen aus Gottes Befehl. Derselbige Befehl aber gehet uns im Neuen Testament nichts an, da sollen wir bleiben bei dem Gelübde, das wir in der Taufe gethan haben: wie wir dem Teufel und allen seinen Werken entgegen sein, und uns an Gottes Güte und Gnade, so er in seinem Sohn Christo Jesu uns bewiesen, halten wollen. Bei dem Gelübde sollen wirs lassen bleiben, und haben unser Leben lang genug damit zu thun, dürfen uns nach andern nicht umsehen. Denn es bringt mit sich öffentliches Bekenntniß des Glaubens, Kreuz, rechte Anrufung im Kreuz, Hoffnung, Glaube und dergleichen anderes.

Darum, wenn die Papisten mit der Schrift einher prallen, es stehe in Mose, im Psalter, in Propheten, was man Gott gelobe, soll man halten; so antworte du und sprich: Ja es ist wahr, wenn mans verstehet vom Gelübde in der Taufe, denn das ist ein ewiges Gelübde. Aber man besehe Mosen und die Propheten fleißig, so wird man finden, daß sie von ihren Gelübden, so im Gesetze gebräuchlich, reden; dieselbigen Gelübde gehen uns nichts an. Darum binden sie uns eben so wenig, als andere Sprüche oder Gebote, die da verbieten, daß man Schweinen-Fleisch, Hasen, Aale und dergleichen nicht essen soll.

Eben also ists eine große Thorheit, daß die Papisten mit den Sprüchen, so auf jene Gelübde gerichtet, ihre Gelübde gedenken zu erhalten, von welchen Gott nie keinen Befehl ihnen gegeben hat, ja welche wider Gottes Wort und Ordnung und derhalben ärgerlich und verboten sind.

Von bürgerlichem Gelübden und wie fern dieselben zu halten, handelt man hie nicht, sondern von den Gelübden allein, die da sollen ein Gottesdienst sein. Da haben wir im Neuen Testament keines, denn das Gelübde von der Taufe. Zu andern guten Werken dürfen wir uns nicht geloben, wir sinds vorhin zu halten schuldig, sintemal Gott sie uns allen hat auferlegt und zu halten geboten. Das sei genug zum Bericht von den Gelübden.
